\SourcePage{87}
\SGAchapter{IX}{代数几何与形式几何}{Géométrie algébrique et géométrie formelle}
本讲旨在把 EGA III 的定理 3.4.2 和 4.1.5 推广到\Term{morphism}{态射}不\Term{proper}{固有}的情形。
\SGAsection{1}{比较定理}
设 $f:X\to X'$ 是\Term{separated}{分离}的\Term{finite-type}{有限型}\Term{prescheme}{预概形}\Term{morphism}{态射}，并假设 $X'$ \Term{locally-noetherian}{局部诺特}。设 $Y'$ 是 $X'$ 的\Term{closed-subset}{闭子集}，$Y=f^{-1}(Y')$。以 $\widehat X$、$\widehat{X'}$ 分别表示 $X$、$X'$ 沿 $Y$、$Y'$ 的\Term{formal-completion}{形式完备化}，$\widehat f$ 为 $f$ \Term{completion}{完备化}后所得的\Term{morphism}{态射}。
\[\begin{tikzcd}X\arrow[d,"f" ]&Y\arrow[l]\arrow[d]\\X'&Y'\arrow[l]\end{tikzcd}\qquad
\begin{tikzcd}\widehat X\arrow[r,"j"]\arrow[d,"\widehat f" ]&X\arrow[d,"f"]\\\widehat{X'}\arrow[r,"i"]&X'\end{tikzcd}\tag{1.1}\]
以 $j$（分别 $i$）表示从 $\widehat X$ 到 $X$（分别从 $\widehat{X'}$ 到 $X'$）的\Term{homomorphism}{同态}。已知 $i,j$ \Term{flat}{平坦}。设 $\mathcal I'$ 是 $Y'$ 的\Term{ideal-of-definition}{定义理想}，令 $\mathcal J=f^*(\mathcal I')\mathcal O_X$，它是 $Y$ 的\Term{ideal-of-definition}{定义理想}。因此
\[\widehat{X'}=(Y',\varprojlim_{k\in\mathbb N}\mathcal O_{X'}/\mathcal I'^{k+1}),\qquad
\widehat X=(Y,\varprojlim_{k\in\mathbb N}\mathcal O_X/\mathcal J^{k+1}).\tag{1.2}\]
对每个 $k\in\mathbb N$，置
\[Y'_k=(Y',\mathcal O_{X'}/\mathcal I'^{k+1}),\qquad Y_k=(Y,\mathcal O_X/\mathcal J^{k+1}).\tag{1.3}\]
设 $F$ 是\Term{coherent}{凝聚} $\mathcal O_X$-模。对每个 $k\in\mathbb N$，置
\[F_k=F/\mathcal J^{k+1}F,\qquad \widehat F=j^*(F)\simeq\varprojlim F_k.\tag{1.4}\]
若置
\[(R^if_*(F))^\wedge=\varprojlim_{k\in\mathbb N}(R^if_*(F)\otimes_{\mathcal O_{X'}}\mathcal O_{Y'_k}),\quad i\in\mathbb Z,\tag{1.5}\]
\SourcePage{88}
则有自然\Term{homomorphism}{同态}
\[r_i:i^*(R^if_*(F))\longrightarrow (R^if_*(F))^\wedge,\tag{1.6}\]
当 $R^if_*(F)$ \Term{coherent}{凝聚}时它是\Term{isomorphism}{同构}。如 EGA III 4.1.1 所述，有\Term{commutative-diagram}{交换图}
\[\begin{tikzcd}
i^*(R^if_*(F))\arrow[r,"\rho_i"]\arrow[d,"r_i" ]&R^i\widehat f_*(\widehat F)\arrow[d,"\psi_i"]\\
(R^if_*(F))^\wedge\arrow[r,"\varphi_i"]&\varprojlim_{k\in\mathbb N}R^if_*(F_k).
\end{tikzcd}\tag{1.7}\]
在所引文献中，已知 $R^if_*(F)$ \Term{coherent}{凝聚}，因而把 (1.6) 的源和靶等同，得到一个\Term{commutative-diagram}{交换图}。在这里，$R^if_*(F)$ 仅对某些 $i$ \Term{coherent}{凝聚}；我们将对这些 $i$ 研究 (1.7)。考虑\Term{graded-ring}{分次环}
\[\mathcal S=\bigoplus_{k\in\mathbb N}\mathcal I'^k,\tag{1.8}\]
以及\Term{graded}{分次} $\mathcal S$-模
\[\mathcal H^i=\bigoplus_{k\in\mathbb N}R^if_*(\mathcal J^kF),\quad i\in\mathbb Z,\tag{1.9}\]
其 $\mathcal S$-模结构如下定义。\Term{sheaf}{层} $R^if_*(\mathcal J^kF)$ 由\Term{presheaf}{预层}\Term{adjunction}{伴随}而来，后者给 $X'$ 的每个\Term{affine}{仿射}\Term{open-subset}{开集} $U'$ 配上
\[H^i(f^{-1}(U'),\mathcal J^kF|_{f^{-1}(U')}).\tag{1.10}\]
因此设 $U'$ 是 $X'$ 的\Term{affine}{仿射}\Term{open-subset}{开集}，置 $U=f^{-1}(U')$。设 $x'\in\mathcal I'^m(U')$，并以 $x$ 表示它在 $\mathcal J^m(U)$ 中的像。$F|_U$ 上的乘 $x$ 映射把 $\mathcal J^kF|_U$ 映到 $\mathcal J^{k+m}F|_U$，由\Term{functor}{函子}性得到\Term{morphism}{态射}
\[\mu^i_{x',k}(U'):H^i(U,\mathcal J^kF|_U)\longrightarrow H^i(U,\mathcal J^{k+m}F|_U),\tag{1.11}\]
对所有 $i\in\mathbb Z,k\in\mathbb N$ 定义；取\Term{adjunction}{伴随}\Term{sheaf}{层}便给出 $\mathcal H^i$ 的\Term{graded}{分次} $\mathcal S$-模结构。
\begin{SGAresult}{定理}{1.1}
设 $n$ 为整数。假设\Term{graded}{分次} $\mathcal S$-模 $\mathcal H^i$ 在 $i=n-1,n$ 时\Term{finite-type}{有限型}。则：
\begin{enumerate}
\item[(0)] $r_n,r_{n-1}$ 是\Term{isomorphism}{同构}，且 $R^i\widehat f_*(\widehat F)$ 在 $i=n-1$ 时\Term{coherent}{凝聚}；
\item[(1)] 对 $i=n-1$，$\rho_i,\varphi_i,\psi_i$ 是\Term{topology}{拓扑}\Term{isomorphism}{同构}。特别地，$R^{n-1}f_*(F)$ 上由\Term{homomorphism}{同态}
\[R^{n-1}f_*(F)\longrightarrow R^{n-1}f_*(F_k)\tag{1.12}\]
的核定义的\Term{filtration}{滤过}是 $\mathcal I'$-良\Term{filtration}{滤过}；
\SourcePage{89}
\item[(2)] 对 $i=n$，$\rho_i,\varphi_i,\psi_i$ 是\Term{monomorphism}{单态射}；此外 $R^nf_*(F)$ 上的\Term{filtration}{滤过}是 $\mathcal I'$-良\Term{filtration}{滤过}，且 $\psi_n$ 是\Term{isomorphism}{同构}；
\item[(3)] 对 $i=n-2,n-1$，\Term{projective-system}{投射系} $(R^if_*(F_k))$ 满足一致 Mittag--Leffler 条件，即存在固定整数 $k\geq0$，使对每个 $p\geq0$ 及 $p'\geq p+k$，有
\[\operatorname{Im}[R^if_*(F_{p'})\to R^if_*(F_p)]=\operatorname{Im}[R^if_*(F_{p+k})\to R^if_*(F_p)].\]
\end{enumerate}
\end{SGAresult}
\begin{proof}
仿照 EGA III 4.1.8，容易归结到 $X'$ 是\Term{noetherian}{诺特}环 $A$ 的谱的情形。这时已知
\[R^if_*(F)=\widetilde{H^i(X,F)}\qquad\text{（参见 1.10）}.\tag{1.13}\]
设 $I$ 是满足 $\widetilde I=\mathcal I'$ 的 $A$ 的\Term{ideal}{理想}，并置
\[S=\bigoplus_{k\in\mathbb N}I^k,\tag{1.14}\]
\[H^i=\bigoplus_{k\in\mathbb N}H^i(X,\mathcal J^kF),\qquad i\in\mathbb Z,\tag{1.15}\]
其中 $H^i$ 带有取 $U'=X'$ 时由 1.11 定义的\Term{graded}{分次} $S$-模结构。证明仿照 EGA III 4.1.5，现概述如下。我们考察 $\varphi_i,\psi_i$，对应于模\Term{homomorphism}{同态}
\[\begin{tikzcd}&H^i(\widehat X,\widehat F)\arrow[d,"\psi_i"]\\H^i(X,F)^\wedge\arrow[r,"\varphi_i"]&\varprojlim_k H^i(X,F_k).
\end{tikzcd}\tag{1.16}\]
(a) 仅假设 $H^i$ 是\Term{finite-type}{有限型}\Term{graded}{分次} $S$-模。由此推出 $H^i(X,F)$ 上由模
\[R^i_k=\ker(H^i(X,F)\to H^i(X,F_k))\tag{1.17}\]
定义的\Term{filtration}{滤过}为 $I$-良\Term{filtration}{滤过}。为此利用\Term{cohomology}{上同调}\Term{exact-sequence}{正合列}
\[H^i(X,\mathcal J^{k+1}F)\longrightarrow H^i(X,F)\longrightarrow H^i(X,F_k),\tag{1.18}\]
它表明\Term{graded}{分次} $S$-模 $\bigoplus_{k\in\mathbb N}R^i_k$ 是 $H^i$ 的\Term{graded}{分次}子 $S$-模 $\bigoplus_{k\in\mathbb N}H^i(X,\mathcal J^{k+1}F)$ 的商，因此\Term{finite-type}{有限型}，因为 $S$ \Term{noetherian}{诺特}。第一点由此得证。置
\[M^i=H^i(X,F),\qquad H^i_k=H^i(X,F_k).\tag{1.19}\]
\SourcePage{90}
有\Term{commutative-diagram}{交换图}
\[\begin{tikzcd}H^i(X,F)^\wedge\arrow[r,"s_i"]\arrow[dr,"\varphi_i"']&\varprojlim_k(M^i/R^i_k)\arrow[d,"t_i"]\\&\varprojlim_k H^i_k,\end{tikzcd}\tag{1.20}\]
其中 $s_i$ 是\Term{isomorphism}{同构}，因为 $H^i(X,F)$ 的\Term{filtration}{滤过}为 $I$-良\Term{filtration}{滤过}。此外 $t_i$ 是\Term{monomorphism}{单态射}：\Term{functor}{函子} $\varprojlim$ \Term{left-exact}{左正合}，而根据 $R^i_k$ 的定义，对每个 $k\geq0$，自然\Term{morphism}{态射} $M^i/R^i_k\to H^i_k$ 都是\Term{monomorphism}{单态射}。
为研究 $t_i$ 的\Term{surjectivity}{满射}性，引入
\[Q^i_k=\operatorname{coker}(H^i(X,F)\to H^i(X,F_k)),\tag{1.21}\]
从而有\Term{exact-sequence}{正合列}的\Term{projective-system}{投射系}
\[0\longrightarrow R^i_k\longrightarrow M^i\longrightarrow H^i_k\longrightarrow Q^i_k\longrightarrow0.\tag{1.22}\]
利用\Term{cohomology}{上同调}\Term{exact-sequence}{正合列}
\[H^i(X,F)\longrightarrow H^i(X,F_k)\longrightarrow H^{i+1}(X,\mathcal J^{k+1}F),\tag{1.23}\]
可见\Term{graded}{分次} $S$-模
\[Q^i=\bigoplus_{k\in\mathbb N}Q^i_k\tag{1.24}\]
是 $H^{i+1}$ 的\Term{graded}{分次}子 $S$-模。此外，对每个 $k\geq0$，有
\[I^{k+1}Q^i_k=0,\tag{1.25}\]
因为 $Q^i_k$ 是 $H^i_k$ 的像。
(b) 仅假设 $H^{i+1}$ \Term{finite-type}{有限型}，并考察 $t_i$（暂且不顾 $s_i$）。因为 $S$ \Term{noetherian}{诺特}，$Q^i$ \Term{finite-type}{有限型}；因 $I^{k+1}Q^i_k=0$，可得存在整数 $r\geq0,k_0\geq0$，使
\[I^rQ^i_k=0\quad(k\geq k_0).\tag{1.26}\]
由此\Term{projective-system}{投射系} $(Q^i_k)_{k\in\mathbb N}$ \Term{essentially-zero}{本质为零}，因而\Term{projective-system}{投射系} $(H^i_k)_{k\in\mathbb N}$ 满足一致 Mittag--Leffler 条件。由\Term{exact-sequence}{正合列} (1.22) 得到\Term{exact-sequence}{正合列}
\[0\longrightarrow M^i/R^i_k\longrightarrow H^i_k\longrightarrow Q^i_k\longrightarrow0,\tag{1.27}\]
继而得到\Term{exact-sequence}{正合列}
\[0\longrightarrow\varprojlim_k M^i/R^i_k\xrightarrow{t_i}\varprojlim_k H^i_k\longrightarrow\varprojlim_k Q^i_k.\tag{1.28}\]
然而 $(Q^i_k)$ \Term{essentially-zero}{本质为零}，所以 $t_i$ 是\Term{isomorphism}{同构}。
(c) 证明若 $H^i$ \Term{finite-type}{有限型}，则 $\psi_i$ 是\Term{isomorphism}{同构}。只须以 $X$ 的\Term{affine}{仿射}\Term{open-subset}{开集}为\Term{open-subset}{开集}基应用 EGA $0_{\mathrm{III}}$ 13.3.1。这是允许的，
\SourcePage{91}
因为由 (b)，\Term{projective-system}{投射系} $(H^{i-1}_k)_{k\in\mathbb N}$ 满足 Mittag--Leffler 条件。定理由 (a)、(b)、(c) 形式地推出。注意，实际上在证明的每一步，仅对一个 $i$ 用到了 $H^i$ 的\Term{finite}{有限}性。
\end{proof}
下面给出几个满足定理 1.1 假设的例子。
\begin{SGAresult}{推论}{1.2}
假设 $\mathcal I'$ 由 $\mathcal O_{X'}$ 的一个\Term{section}{截面} $t'$ 生成，以 $t$ 表示 $\mathcal O_X$ 的对应\Term{section}{截面}。设 $F$ 是\Term{coherent}{凝聚} $\mathcal O_X$-模，$n$ 是整数。假设：
(i) $t$ 为 $F$-\Term{regular}{正则}元（即 $F$ 上的乘 $t$ 映射为\Term{monomorphism}{单态射}）；
(ii) $R^if_*(F)$ 在 $i=n-1,n$ 时\Term{coherent}{凝聚}。
则定理 1.1 的假设成立。
\end{SGAresult}
\begin{proof}
注意乘 $t^k$ 给出\Term{isomorphism}{同构} $F\xrightarrow\sim\mathcal J^kF$，从而
\[\mathcal H^i\simeq R^if_*(F)\otimes_{\mathcal O_{X'}}\mathcal O_{X'}[T],\tag{1.29}\]
其中 $T$ 为不定元。结论随即得到。
\end{proof}
\begin{SGAresult}{推论}{1.3}
假设 $X'=\operatorname{Spec}(A)$，其中 $A$ 为对 $I$-\Term{adic-topology}{进拓扑}\Term{separated}{分离}且\Term{complete}{完备}的\Term{noetherian}{诺特}环。假设 $S$-模 $H^i$ 在 $i=n-1,n$ 时\Term{finite-type}{有限型}（参见 1.14、1.15）。则定理 1.1 的假设成立，且有\Term{isomorphism}{同构}的\Term{commutative-diagram}{交换图}\footnote{编者注：同一方向可参阅 Chow W.-L., ``Formal functions on homogeneous spaces'', Invent. Math. 86 (1986), no.~1, pp.~115--130。作者证明如下结果：设 $X$ 为域上的\Term{algebraic-variety}{代数簇}，在\Term{algebraic-group}{代数群} $G$ \Term{action}{作用}下齐性；设 $Z$ 为 $X$ 的正维\Term{complete}{完备}子簇。假设 $Z$ 以下述意义生成 $X$：给定 $p\in Z$，令 $\Gamma_p$ 为把 $p$ 映入 $Z$ 的 $G$ 的元素\Term{set}{集合}。若 $\Gamma_p$ 的含 $1$ 的\Term{connected}{连通}分支生成的群是整个 $G$，就说 $Z$ 生成。在此情形，$X$ 沿 $Z$ 的每个形式有理函数都是\Term{algebra}{代数}的；可与原书第 138 页编者注 (3) 引述的 Hironaka 和 Matsumura 的结果比较。沿这些作者引入的技术，值得指出 Gieseker 的一个很漂亮的\Term{algebraization}{代数化}结果：Gieseker D., ``On two theorems of Griffiths about embeddings with ample normal bundle'', Amer. J. Math. 99 (1977), no.~6, pp.~1137--1150，定理 4.1、4.2。设 $X$ 是\Term{algebraically-closed-field}{代数闭域}上正维\Term{connected}{连通}\Term{projective-variety}{射影簇}，\Term{local-scheme}{局部}为\Term{complete-intersection}{完全交}。假设 $X$ 有两个嵌入，分别嵌入\Term{smooth}{光滑}\Term{projective-variety}{射影簇} $Y,W$。若 $X$ 在 $Y,W$ 中的\Term{formal-completion}{形式完备化}\Term{equivalence}{等价}，则存在含 $X$ 为闭子\Term{scheme}{概形}的\Term{scheme}{概形} $U$，它可嵌入 $Y,W$，成为 $X$ 在 $Y,W$ 中的\Term{etale}{平展}邻域。换言之，形式\Term{equivalence}{等价}蕴涵\Term{etale}{平展}\Term{equivalence}{等价}。另见 Faltings G., ``Formale Geometrie and homogene Raüme'', Invent. Math. 64 (1981), pp.~123--165。}
\[\begin{tikzcd}H^i(X,F)\arrow[r,"\rho'_i"]\arrow[dr,"\varphi'_i"']&H^i(\widehat X,\widehat F)\arrow[d,"\psi_i"]\\&\varprojlim_kH^i(X,F_k)\end{tikzcd}\qquad(i=n-1).\tag{1.1}\]
\end{SGAresult}
\SourcePage{92}
\begin{proof}
只须注意 $H^i(X,F)$ \Term{finite-type}{有限型}，故\Term{isomorphism}{同构}于其\Term{completion}{完备化}。将模\Term{sheaf}{层}的图 (1.7) 改写到 $A$-模\Term{category}{范畴}中，并把左侧竖\Term{arrow}{箭头}换为 $H^i(X,F)$，就得到 (1.1)。
\end{proof}
\begin{SGAresult}{命题}{1.4}
设 $A$ 是\Term{noetherian}{诺特}环，$t\in A$，并假设 $A$ 对 $(tA)$-\Term{adic-topology}{进拓扑}\Term{separated}{分离}且\Term{complete}{完备}。置
\[X'=\operatorname{Spec}(A),\quad Y'=V(t),\quad I=(tA).\tag{1.31}\]
设 $T$ 为 $X'$ 的\Term{closed-subset}{闭子集}，置
\[X=X'-T,\quad Y=Y'\cap X=Y'-(Y'\cap T).\tag{1.32}\]
设 $F$ 是\Term{coherent}{凝聚} $\mathcal O_X$-模。最后置
\[T'=\{x\in X'\mid\operatorname{codim}(\overline{\{x\}}\cap T,\overline{\{x\}})=1\}.\tag{1.33}\]
假设：a) $t$ 为 $F$-\Term{regular}{正则}元；b) $\operatorname{prof}_{T'}(F)\geq n+1$；c) $A$ 为\Term{regular}{正则}\Term{noetherian}{诺特}环的商。
则在图 (1.1) 中，\Term{morphism}{态射} $\rho'_i,\varphi'_i,\psi_i$ 在 $i<n$ 时为\Term{isomorphism}{同构}，在 $i=n$ 时为\Term{monomorphism}{单态射}。此外 $\psi_n$ 是\Term{isomorphism}{同构}。
\end{SGAresult}
\begin{proof}
由 1.3、1.2，只须证明 $R^if_*(F)$ 在 $i\leq n$ 时\Term{coherent}{凝聚}；这由\Term{finiteness-theorem}{有限性定理} 2.1 得出。
\end{proof}
特别地：
\begin{SGAresult}{例}{1.5}
当 $A$ 是\Term{local-ring}{局部环}且 $t$ 属于 $A$ 的根 $\mathfrak r(A)$ 时，我们将应用 1.4。这时取 $T=\{\mathfrak r(A)\}$。在 $n=1$ 时得到如下陈述：
若\Term{noetherian}{诺特}环 $A$ 对 $t$-\Term{adic-topology}{进拓扑}\Term{separated}{分离}且\Term{complete}{完备}，且为\Term{regular}{正则}环的商（例如 $A$ \Term{complete}{完备}），若此外 $t$ 为 $F$-\Term{regular}{正则}元，且对每个满足 $\dim A/x=1$ 的 $x\in\operatorname{Spec}(A)$ 有 $\operatorname{prof}F_x\geq2$，则自然\Term{homomorphism}{同态}
\[\Gamma(X,F)\longrightarrow\Gamma(\widehat X,\widehat F)\]
为\Term{isomorphism}{同构}。
\end{SGAresult}
事实上，保持 1.4 的记号，有 $T=\{\mathfrak r(A)\}$，而公式 (1.33) 给出
\[T'=\{x\in\operatorname{Spec}(A)\mid\dim A/x=1\}.\]
\SourcePage{93}
\SGAsection{2}{存在定理}
先把 EGA III 3.4.2 表述为稍一般的形式。设 $f:X\to X'$ 是\Term{formal-prescheme}{形式预概形}的\Term{adic-morphism}{进态射}\footnote{这个假设并非本质的，参见第 XII 讲，原书第 118 页。}，且 $X'$ \Term{noetherian}{诺特}。设 $\mathcal I'$ 是 $X'$ 的\Term{ideal-of-definition}{定义理想}；因为 $f$ 为\Term{adic-morphism}{进态射}，$f^*\mathcal I'=\mathcal J$ 是 $X$ 的\Term{ideal-of-definition}{定义理想}\footnote{编者注：这正是定义，参见 EGA I 10.12.1。}。对每个 $n\in\mathbb N$，置
\[X_n=(X,\mathcal O_X/\mathcal J^{n+1}),\tag{2.1}\]
它是与 $X$ 有同一下\Term{sheaf}{层}\Term{topological-space}{拓扑空间}的通常\Term{prescheme}{预概形}。设 $F$ 是\Term{coherent}{凝聚} $\mathcal O_X$-模。对每个 $k\in\mathbb N$，$\mathcal O_{X_k}$-模
\[F_k=F/\mathcal J^{k+1}F\tag{2.2}\]
\Term{coherent}{凝聚}。对每个 $i$ 有\Term{homomorphism}{同态}
\[\psi_i:R^if_*(F)\longrightarrow\varprojlim_kR^if_*(F_k),\tag{2.3}\]
由自然\Term{homomorphism}{同态} $F\to F_k$\hfill(2.4) 通过\Term{functor}{函子}性得到。置
\[\mathcal S=\operatorname{gr}_{\mathcal I'}\mathcal O_{X'}=\bigoplus_{k\in\mathbb N}\mathcal I'^k/\mathcal I'^{k+1},\tag{2.5}\]
\[\operatorname{gr}_{\mathcal J}(F)=\bigoplus_{k\in\mathbb N}\mathcal J^kF/\mathcal J^{k+1}F,\tag{2.6}\]
\[K^i=R^if_*(\operatorname{gr}_{\mathcal J}(F))=\bigoplus_{k\in\mathbb N}R^if_*(\mathcal J^kF/\mathcal J^{k+1}F).\tag{2.7}\]
显然 $K^i$ 带有\Term{graded}{分次} $\mathcal S$-模结构。
\begin{SGAresult}{定理}{2.1}
假设 $K^i$ 在 $i=n-1,n,n+1$ 时为\Term{finite-type}{有限型}\Term{graded}{分次} $\mathcal S$-模。则：
(i) $R^nf_*(F)$ \Term{coherent}{凝聚}；
(ii) \Term{homomorphism}{同态} $\psi_n$ (2.3) 为\Term{topology}{拓扑}\Term{isomorphism}{同构}，且 (2.3) 右侧的自然\Term{filtration}{滤过}为 $\mathcal I'$-良\Term{filtration}{滤过}；
(iii) \Term{projective-system}{投射系} $(R^nf_*(F_k))$ 满足一致 Mittag--Leffler 条件。
\end{SGAresult}
证明由 EGA $0_{\mathrm{III}}$ 13.7.7（参见 EGA III 3.4.2）很容易得出，但须按 EGA III 2, ErrIII 24 所指示的方式修正其第 78 页的文字。
\SourcePage{94}
\begin{SGAresult}{定理}{2.2}
设 $A$ 是\Term{noetherian}{诺特}\Term{adic-ring}{进环}，$I$ 是 $A$ 的\Term{ideal-of-definition}{定义理想}。设 $T$ 是 $X'=\operatorname{Spec}(A)$ 的\Term{closed-subset}{闭子集}。假设 $I$ 由某个 $t\in A$ 生成。沿用 1.31、1.32、1.33 的记号。设 $\mathfrak F$ 是\Term{coherent}{凝聚} $\mathcal O_{\widehat X}$-模。置
\[F_0=\mathfrak F/\mathcal J\mathfrak F,\tag{2.1}\]
其中 $\mathcal J=t\mathcal O_{\widehat X}$ 是 $\widehat X$ 的\Term{ideal-of-definition}{定义理想}。假设 $A$ 为\Term{regular}{正则}\Term{noetherian}{诺特}环的商，且：
(1) $t$ 为 $\mathfrak F$-\Term{regular}{正则}元；(2) $\operatorname{prof}_{T'}F_0\geq2$。
则存在\Term{coherent}{凝聚} $\mathcal O_X$-模 $F$，使 $\widehat F\simeq\mathfrak F$。\footnote{编者注：后来已得到许多\Term{algebraization}{代数化}陈述。除下文所引者外，参见原书第 155 页编者注 (22) 和第 203 页编者注 (7) 所引 Faltings、Raynaud 夫人的论文。特别想到 Artin 的结果，尤其 Artin M., ``Algebraization of formal moduli. I'', in Global Analysis (Papers in Honor of K. Kodaira), Univ. Tokyo Press, Tokyo, 1969, pp.~21--71；以及较近期关于叶状结构的叶之\Term{algebra}{代数}性的结果，特别见 Bost J.-B., ``Algebraic leaves of algebraic foliations over number fields'', Publ. Math. Inst. Hautes Études Sci. 93 (2001), pp.~161--221；Chambert-Loir A., ``Théorèmes d'algébricité en géométrie diophantienne (d'après J.-B. Bost, Y. André, D. \& G. Chudnovsky)'', in Séminaire Bourbaki, Vol.~2000/2001, Astérisque, vol.~282, Société mathématique de France, Paris, 2002, Exp.~886, pp.~175--209 及其中参考文献。尤其在这两篇文章中，可找到\Term{algebraization}{代数化}问题与丢番图逼近理论之间联系的讨论。}
\end{SGAresult}
\begin{proof}
只须证明 $\widehat f_*(\mathfrak F)$ 是\Term{coherent}{凝聚} $\mathcal O_{\widehat X}$-模，其中 $\widehat f:\widehat X\to\widehat{X'}$ 是由 $X\hookrightarrow X'$ 作 $t$-进\Term{completion}{完备化}所得的\Term{formal-prescheme}{形式预概形}\Term{morphism}{态射}。事实上，$A$ 对 $t$-\Term{adic-topology}{进拓扑}\Term{separated}{分离}且\Term{complete}{完备}，因此存在 $A$-模 $F'$，其\Term{completion}{完备化}\Term{isomorphism}{同构}于 $\widehat f_*(\mathfrak F)$。因为 $X$ 是 $X'$ 的\Term{open-subset}{开集}，可取 $F=\widetilde{F'}|_X$。
还须证明可对\Term{formal-prescheme}{形式预概形}\Term{morphism}{态射} $\widehat f$ 及 $\mathfrak F$ 应用 2.1。由假设 (1)，对每个 $k\in\mathbb N$ 有\Term{isomorphism}{同构}
\[\mathcal J^k\mathfrak F/\mathcal J^{k+1}\mathfrak F\longrightarrow\mathfrak F/\mathcal J\mathfrak F,\]
因此只要知道 $R^if_*(F_0)$ 在 $i\leq1$ 时\Term{coherent}{凝聚}，2.1 的假设就成立。这由 (2) 和\Term{finiteness-theorem}{有限性定理} 2.1 得出。结论得证。
\end{proof}
现在只须在 $A$ 为\Term{local-ring}{局部环}的假设下把此陈述\Term{specialization}{特殊化}。
\begin{SGAresult}{推论}{2.3}
设 $A$ 是\Term{noetherian}{诺特}\Term{local-ring}{局部环}，$t\in\mathfrak r(A)$ 为其根中的元素。假设 $A$ 对 $t$-\Term{adic-topology}{进拓扑}\Term{separated}{分离}且\Term{complete}{完备}，且为\Term{regular}{正则}环的商（例如假设 $A$ 是\Term{complete}{完备}\Term{noetherian}{诺特}\Term{local-ring}{局部环}）。置
\[X'=\operatorname{Spec}A,\qquad T=\{\mathfrak r(A)\},\tag{2.9}\]
\SourcePage{95}
并沿用 (1.31)、(1.32)、(1.33) 的记号。设 $\mathfrak F$ 是 $\mathcal O_{\widehat X}$-模。假设：
(1) $t$ 为 $\mathfrak F$-\Term{regular}{正则}元；
(2) $\operatorname{prof}_{T'}F_0\geq2$，其中 $F_0=\mathfrak F/\mathcal J\mathfrak F$，$\mathcal J=t\mathcal O_{\widehat X}$。
则存在\Term{coherent}{凝聚} $\mathcal O_X$-模 $F$，使 $\widehat F\simeq\mathfrak F$。
\end{SGAresult}
注意这里 $T'$ 是 $A$ 中满足 $\dim A/\mathfrak p=1$ 的\Term{prime-ideal}{素理想} $\mathfrak p$ 的\Term{set}{集合}。
\SourcePage{96}
\SourcePage{97}
\SGAchapter{X}{对基本群的应用}{Application au groupe fondamental}
本讲始终以 $X$ 表示\Term{locally-noetherian}{局部诺特}\Term{prescheme}{预概形}，$Y$ 表示 $X$ 的\Term{closed-subset}{闭子集}，$U$ 表示 $Y$ 在 $X$ 中的可变\Term{open-neighborhood}{开邻域}，$\widehat X$ 表示 $X$ 沿 $Y$ 的\Term{formal-completion}{形式完备化}（EGA I 10.8）。对任意\Term{prescheme}{预概形} $Z$，以 $\operatorname{Et}(Z)$ 表示 $Z$ 的\Term{etale}{平展}\Term{covering}{覆盖}\Term{category}{范畴}，以 $L(Z)$ 表示 $Z$ 上\Term{locally-free}{局部自由}\Term{coherent}{凝聚}模的\Term{category}{范畴}。
\SGAsection{1}{比较 Et(完备化) 与 Et(Y)}
设 $\mathcal I$ 是 $Y$ 在 $X$ 中的\Term{ideal-of-definition}{定义理想}。对每个 $n\in\mathbb N$，置 $Y_n=(Y,(\mathcal O_X/\mathcal I^{n+1})|_Y)$。$Y_n$ 构成通常\Term{prescheme}{预概形}的\Term{inductive-system}{归纳系}；给结构\Term{sheaf}{层}赋予\Term{discrete}{离散}\Term{topology}{拓扑}后，也可看作\Term{formal-prescheme}{形式预概形}的\Term{inductive-system}{归纳系}。已知（EGA I 10.6.2），$\widehat X$ 在\Term{formal-prescheme}{形式预概形}\Term{category}{范畴}中是 $Y_n$ 的\Term{direct-limit}{归纳极限}。又已知（EGA I 10.13），给定\Term{finite-type}{有限型} $\widehat X$-\Term{formal-prescheme}{形式预概形} $R$，\Term{equivalence}{等价}于给定\Term{finite-type}{有限型}通常 $Y_n$-\Term{prescheme}{预概形} $R_n$ 的\Term{inductive-system}{归纳系}，满足 $R_n\simeq R_{n+1}\times_{Y_{n+1}}Y_n$。此外，$R$ 是 $\widehat X$ 的\Term{etale}{平展}\Term{covering}{覆盖}，当且仅当每个 $R_n$ 都是 $Y_n$ 的\Term{etale}{平展}\Term{covering}{覆盖}。容易看到\Term{nilpotent}{幂零}元不影响\Term{etale}{平展}\Term{covering}{覆盖}（SGA 1 8.3），即对每个 $n\in\mathbb N$，\Term{base-change}{基变换}\Term{functor}{函子}
\[\operatorname{Et}(Y_{n+1})\longrightarrow\operatorname{Et}(Y_n)\]
是\Term{equivalence-of-categories}{范畴等价}。因此：
\begin{SGAresult}{命题}{1.1}
沿用上述记号，自然\Term{functor}{函子} $\operatorname{Et}(\widehat X)\to\operatorname{Et}(Y)$ 是\Term{equivalence-of-categories}{范畴等价}（参见 SGA 1 8.4）。
\end{SGAresult}
\SGAsection{2}{比较 Et(Y) 与可变 U 的 Et(U)}
引入两个条件，预告的\Term{comparison-theorem}{比较定理}将由此容易推出。设 $X$ 为\Term{locally-noetherian}{局部诺特}\Term{prescheme}{预概形}，$Y$ 为 $X$ 的\Term{closed-subset}{闭子集}。称对 $(X,Y)$ 满足 \Term{lefschetz-condition}{Lefschetz 条件}，记作 $\operatorname{Lef}(X,Y)$，如果对 $X$ 中每个包含 $Y$ 的\Term{open-subset}{开集} $U$ 以及 $U$ 上每个\Term{locally-free}{局部自由}\Term{coherent}{凝聚}\Term{sheaf}{层} $E$，自然\Term{homomorphism}{同态}
\SourcePage{98}
\[\Gamma(U,E)\longrightarrow\Gamma(\widehat X,\widehat E)\]
为\Term{isomorphism}{同构}。
称对 $(X,Y)$ 满足\Term{effective-lefschetz-condition}{有效 Lefschetz 条件}，记作 $\operatorname{Leff}(X,Y)$，如果满足 $\operatorname{Lef}(X,Y)$，并且对 $\widehat X$ 上每个\Term{locally-free}{局部自由}\Term{coherent}{凝聚}\Term{sheaf}{层} $\mathcal E$，存在 $Y$ 的\Term{open-neighborhood}{开邻域} $U$、$U$ 上的\Term{locally-free}{局部自由}\Term{coherent}{凝聚}\Term{sheaf}{层} $E$，以及\Term{isomorphism}{同构} $\widehat E\simeq\mathcal E$。
这些条件在两个重要例子中成立。
\begin{SGAresult}{例}{2.1}
设 $A$ 为\Term{noetherian}{诺特}环，$t\in\mathfrak r(A)$ 为属于 $A$ 的根的 $A$-\Term{regular}{正则}元。假设 $A$ 为\Term{regular}{正则}\Term{local-ring}{局部环}的商，并对 $t$-\Term{adic-topology}{进拓扑}\Term{complete}{完备}（例如 $A$ 对 $\mathfrak r(A)$-\Term{adic-topology}{进拓扑}\Term{complete}{完备}）。置 $X'=\operatorname{Spec}(A)$、$Y'=V(t)$；另置 $x=\mathfrak r(A)$、$X=X'-\{x\}$、$Y=Y'-\{x\}$。于是 $X$ 在 $X'$ 中开，且 $Y=X\cap Y'$。则：
(i) 若对 $A$ 的每个满足 $\dim A/\mathfrak p=1$ 的\Term{prime-ideal}{素理想} $\mathfrak p$（即 $X$ 的每个\Term{closed-point}{闭点}）有 $\operatorname{prof}A_{\mathfrak p}\geq2$，则有 $\operatorname{Lef}(X,Y)$；
(ii) 若此外，对每个满足 $t\in\mathfrak p$ 且 $\dim A/\mathfrak p=1$ 的\Term{prime-ideal}{素理想} $\mathfrak p$（即 $Y$ 的每个\Term{closed-point}{闭点}）有 $\operatorname{prof}A_{\mathfrak p}\geq3$，则有 $\operatorname{Leff}(X,Y)$。\footnote{编者注：(i) 可以稍作改进，见 Raynaud M., ``Théorèmes de Lefschetz en cohomologie des faisceaux cohérents et en cohomologie étale. Application au groupe fondamental'', Ann. Sci. Éc. Norm. Sup. (4) 7 (1974), pp.~29--52，推论 I.1.4、I.5；条件 (ii) 可改进，从而去掉沿 $Y$ 的\Term{depth}{深度}条件（准确陈述见该文定理 3.3）。后一点的证明非常技术化；上述文章仅给证明提示，并参照先前发表于 Bulletin de la Société mathématique de France 的详细版本。}
\end{SGAresult}
\begin{proof}
先证明 $Y$ 在 $X$ 中每个\Term{open-neighborhood}{开邻域} $U$ 的补集都是\Term{finite}{有限}个（$X$ 中）\Term{closed-point}{闭点}的并。注意 $U$ 在 $X$ 中开，故在 $X'$ 中开，因而 $Z'=X'-U$ 闭。设 $I$ 是 $Z'$ 的\Term{ideal-of-definition}{定义理想}；只须证明 $A/I$ 的\Term{dimension}{维数}为 $1$。因为 $Z'\cap Y'=\{x\}$，所以 $A/(I+(t))$ 为\Term{artinian-ring}{阿廷环}，由主\Term{ideal}{理想}定理即得结论。
第一个假设\Term{equivalence}{等价}于：“对 $A$ 的每个\Term{prime-ideal}{素理想} $\mathfrak p\ne\mathfrak r(A)$，有 $\operatorname{prof}A_{\mathfrak p}\geq3-\dim A/\mathfrak p$。”事实上，$A$ 是\Term{regular}{正则}环的商，因此可对\Term{prescheme}{预概形} $X'$、\Term{closed-subset}{闭子集} $\{x\}$、\Term{coherent}{凝聚}\Term{sheaf}{层} $\mathcal O_{X'}$ 应用第 VIII 讲 2.3；注意对 $p\in U=X'-\{x\}$ 有 $c(p)=\dim(A/\mathfrak p)$，因为 $x$ 是 $X'$ 的\Term{closed-point}{闭点}。
设 $U$ 是 $Y$ 在 $X$ 中的\Term{open-neighborhood}{开邻域}，$E$ 是\Term{locally-free}{局部自由} $\mathcal O_U$-模。置 $Z=X-U$，以 $u:U\to X$ 表示\Term{canonical}{典范}\Term{immersion}{浸入}。先证明 $u_*(E)$ 是\Term{coherent}{凝聚} $\mathcal O_X$-模；\Term{equivalence}{等价}地，若 $E'$ 是 $E$ 到 $X$ 的\Term{coherent}{凝聚}延拓，则 $\mathcal H^i_Z(E')$ 在 $i=0,1$ 时\Term{coherent}{凝聚}。
\SourcePage{99}
为此，对\Term{prescheme}{预概形} $X$、\Term{closed-subset}{闭子集} $Z$、\Term{coherent}{凝聚}\Term{sheaf}{层} $E'$ 应用定理 VIII 2.1。只须检验：对每个满足 $c(p)=1$ 的 $p\in U$，有 $\operatorname{prof}E'_p\geq1$，这里
\[c(p)=\operatorname{codim}(\overline{\{p\}}\cap Z,\overline{\{p\}}).\]
若 $p\in U,c(p)=1$，以 $\mathfrak p$ 表示对应的 $A$ 的\Term{ideal}{理想}，可见 $\dim A/\mathfrak p=2$，因为 $U$ 的补集是\Term{finite}{有限}个\Term{closed-point}{闭点}的并，且 $A$ 为\Term{regular}{正则}环的商。此外 $E$ \Term{locally-free}{局部自由}，故对每个 $p\in\operatorname{Supp}E$ 有 $\operatorname{prof}E_p=\operatorname{prof}\mathcal O_{U,p}$。最后，当 $p\in U,c(p)=1$ 时，
\[\operatorname{prof}E'_p=\operatorname{prof}E_p=\operatorname{prof}\mathcal O_{U,p}=\operatorname{prof}A_{\mathfrak p}\geq3-2=1.\]
现在须证明自然\Term{homomorphism}{同态}
\[\Gamma(U,E)\longrightarrow\Gamma(\widehat X,\widehat E)\tag{1}\]
为\Term{isomorphism}{同构}。置 $\mathfrak E=u_*(E)$，则 $\mathfrak E$ \Term{coherent}{凝聚}，且在 $X$ 的每个\Term{closed-point}{闭点}\Term{depth}{深度}至少为 $2$。由此，$R^if_*(\mathfrak E)$ 在 $i=0,1$ 时\Term{coherent}{凝聚}，其中 $f:X\to X'$ 是 $X$ 到 $X'=\operatorname{Spec}(A)$ 的\Term{canonical}{典范}\Term{immersion}{浸入}（第 VIII 讲）。应用 IX 1.5，由于 $A$ 对 $t$-\Term{adic-topology}{进拓扑}\Term{complete}{完备}，得出
\[\Gamma(U,\mathfrak E)\longrightarrow\Gamma(\widehat X,\widehat{\mathfrak E})\tag{2}\]
为\Term{isomorphism}{同构}。有\Term{commutative-diagram}{交换图}
\[\begin{tikzcd}\Gamma(X,\mathfrak E)\arrow[r,"\simeq"]\arrow[dr,"\simeq"']&\Gamma(U,E)\arrow[d]\\&\Gamma(\widehat X,\widehat E),\end{tikzcd}\]
即得结论。
现在设 $\mathcal E$ 是 $\widehat X$ 上\Term{locally-free}{局部自由}\Term{coherent}{凝聚}\Term{sheaf}{层}。如果已证明 $\mathcal E$ 可\Term{algebraization}{代数化}，即\Term{isomorphism}{同构}于某\Term{coherent}{凝聚} $\mathcal O_X$-模 $E$ 的\Term{formal-completion}{形式完备化}，则容易看出 $E$ 在 $Y$ 的某邻域上\Term{locally-free}{局部自由}，便可证明 $\operatorname{Leff}(X,Y)$。设 $\widehat{X'}$ 是 $A$ 对 $t$-\Term{adic-topology}{进拓扑}的形式谱，它可等同于 $X'$ 沿 $Y'$ 的\Term{formal-completion}{形式完备化}。以 $f$ 表示 $X$ 到 $X'$ 的\Term{canonical}{典范}\Term{immersion}{浸入}，以 $f'$ 表示 $Y$ 到 $Y'$ 的\Term{canonical}{典范}\Term{immersion}{浸入}，以 $\widehat f$ 表示\Term{completion}{完备化}所得的\Term{morphism}{态射}。要使 $\mathcal E$ 可\Term{algebraization}{代数化}，只须 $\widehat f_*(\mathcal E)$ 是\Term{coherent}{凝聚} $\mathcal O_{\widehat X}$-模，因为 $A$ 对 $t$-\Term{adic-topology}{进拓扑}\Term{complete}{完备}。设 $\mathcal I=t\mathcal O_{\widehat X}$，它是 $\widehat X$ 的\Term{ideal-of-definition}{定义理想}。对每个 $n\geq0$，置 $\mathcal E_n=\mathcal E/\mathcal I^{n+1}\mathcal E$。在 $Y$ 的每个\Term{closed-point}{闭点} $y$，$\mathcal E_0$ 的\Term{depth}{深度}至少为 $2$：事实上 $t$ 是 $A$-\Term{regular}{正则}元，故 $\operatorname{prof}\mathcal O_{Y_0,y}=\operatorname{prof}\mathcal O_{X,y}-1\geq2$。所以 $\widehat f_*(\mathcal E)$ \Term{coherent}{凝聚}（IX 2.3）。
\end{proof}
\SourcePage{100}
\begin{SGAresult}{例}{2.2}
（这将用于比较\Term{projective-variety}{射影簇}及其\Term{hyperplane}{超平面}\Term{section}{截面}的\Term{fundamental-group}{基本群}。）
设 $K$ 是域，$X$ 是\Term{proper}{固有} $K$-\Term{prescheme}{预概形}。设 $\mathcal L$ 是\Term{ample}{丰沛}\Term{invertible}{可逆} $\mathcal O_X$-模。设 $t\in\Gamma(X,\mathcal L)$ 是 $\mathcal O_X$-\Term{regular}{正则}元，即对每个\Term{open-subset}{开集} $U$ 及每个\Term{isomorphism}{同构} $u:\mathcal L|_U\to\mathcal O_U$，$u(t)$ 在 $\mathcal O_U$ 中不是\Term{zero-divisor}{零因子}（此条件与 $u$ 无关）。设 $Y=V(t)$ 是 $X$ 中由方程 $t=0$ 定义的子\Term{scheme}{概形}。\footnote{编者注：条件 (ii) 多余，见原书第 90 页的注。} 则：
(i) 若对 $X$ 的每个\Term{closed-point}{闭点} $x$ 有 $\operatorname{prof}\mathcal O_{X,x}\geq2$，则有 $\operatorname{Lef}(X,Y)$；
(ii) 若此外对 $Y$ 的每个\Term{closed-point}{闭点} $y$ 有 $\operatorname{prof}\mathcal O_{X,y}\geq3$，则有 $\operatorname{Leff}(X,Y)$。
\end{SGAresult}
此例将在第 XII 讲详细处理。设 $S$ 为\Term{prescheme}{预概形}。已知（EGA II 6.1.2），把\Term{finite}{有限}\Term{flat}{平坦}\Term{covering}{覆盖} $r:R\to S$ 对应到 $\mathcal O_X$-\Term{algebra}{代数} $r_*(\mathcal O_R)$ 的\Term{functor}{函子}，在 $S$ 的\Term{finite}{有限}\Term{flat}{平坦}\Term{covering}{覆盖}\Term{category}{范畴}与\Term{locally-free}{局部自由}\Term{coherent}{凝聚} $\mathcal O_X$-\Term{algebra}{代数}\Term{category}{范畴}之间诱导\Term{equivalence}{等价}。设 $U$ 为 $Y$ 的\Term{open-neighborhood}{开邻域}，$r:R\to U$ 为 $U$ 的\Term{finite}{有限}\Term{flat}{平坦}\Term{covering}{覆盖}。设 $\widehat R$ 为由\Term{base-change}{基变换}得到的 $\widehat X$ 的\Term{finite}{有限}\Term{flat}{平坦}\Term{covering}{覆盖}。有 $\widehat r_*(\mathcal O_{\widehat R})\simeq\widehat{r_*(\mathcal O_R)}$。
假设 $\operatorname{Lef}(X,Y)$。这意味着对每个 $U$，\Term{inverse-image}{逆像}\Term{functor}{函子} $L(U)\to L(\widehat X)$ \Term{fully-faithful}{全忠实}。事实上，若 $E,F$ 是两个\Term{locally-free}{局部自由}\Term{coherent}{凝聚} $\mathcal O_U$-模，则 $\mathcal Hom(E,F)$ 也\Term{coherent}{凝聚}且\Term{locally-free}{局部自由}；$\mathcal Hom(E,F)$ 也\Term{coherent}{凝聚}且\Term{locally-free}{局部自由}。由假设，自然映射
\[\Gamma_U(\mathcal Hom(E,F))\longrightarrow\Gamma_{\widehat X}(\widehat{\mathcal Hom(E,F)})\]
是\Term{isomorphism}{同构}，结论由此得到，因为在\Term{locally-free}{局部自由}的情况下 $\mathcal Hom$ 与\Term{completion}{完备化}可交换。又\Term{completion}{完备化}与\Term{tensor-product}{张量积}可交换，所以把\Term{locally-free}{局部自由}\Term{coherent}{凝聚} $\mathcal O_U$-\Term{algebra}{代数} $\mathcal A$ 对应到 $\mathcal O_{\widehat X}$-\Term{algebra}{代数} $\widehat{\mathcal A}$ 的\Term{functor}{函子}\Term{fully-faithful}{全忠实}。更进一步，若 $E$ 为\Term{locally-free}{局部自由}\Term{coherent}{凝聚} $\mathcal O_U$-模，则 $\widehat E$ 上的交换 $\mathcal O_{\widehat X}$-\Term{algebraic-structure}{代数结构}之间有一一对应。
\begin{SGAresult}{命题}{2.3}
设 $X$ 为\Term{locally-noetherian}{局部诺特}\Term{prescheme}{预概形}，$Y$ 为 $X$ 的\Term{closed-subset}{闭子集}，$\widehat X$ 为 $X$ 沿 $Y$ 的\Term{formal-completion}{形式完备化}。对每个 $U\supset Y$ 的\Term{open-subset}{开集} $U\subset X$，以 $L_U$（分别 $P_U,E_U$）表示把\Term{locally-free}{局部自由}\Term{coherent}{凝聚} $\mathcal O_U$-模（分别 $U$ 的\Term{finite}{有限}\Term{flat}{平坦}\Term{covering}{覆盖}、$U$ 的\Term{etale}{平展}\Term{covering}{覆盖}）对应到其沿 $\widehat X\to X$ 的\Term{inverse-image}{逆像}的\Term{functor}{函子}。
(i) 若有 $\operatorname{Lef}(X,Y)$，则对 $Y$ 的每个\Term{open-neighborhood}{开邻域} $U$，\Term{functor}{函子} $L_U,P_U,E_U$ \Term{fully-faithful}{全忠实}。
\SourcePage{101}
(ii) 若有 $\operatorname{Leff}(X,Y)$，则对每个\Term{locally-free}{局部自由}\Term{coherent}{凝聚} $\mathcal O_{\widehat X}$-模 $\mathcal E$（分别……），存在\Term{open-subset}{开集} $U$ 及\Term{locally-free}{局部自由}\Term{coherent}{凝聚} $\mathcal O_U$-模 $E$（分别……），使 $L_U(E)\simeq\mathcal E$（分别……）。
\end{SGAresult}
\begin{proof}
(i) 已证。(ii) 对 $L_U,P_U$，由 (i) 和假设得出。此外，若 $R$ 为 $\widehat X$ 的\Term{etale}{平展}\Term{covering}{覆盖}，则存在 $Y$ 在 $X$ 中的\Term{open-neighborhood}{开邻域} $U$ 及 $U$ 的\Term{finite}{有限}\Term{flat}{平坦}\Term{covering}{覆盖} $R'$，使 $\widehat{R'}\simeq R$。所得 $Y$ 的\Term{covering}{覆盖} $R''$ 由 1.1 \Term{etale}{平展}，故 $R'$ 在 $Y$ 的某邻域 $U'$ 上\Term{etale}{平展}。
\end{proof}
\begin{SGAresult}{推论}{2.4}
若有 $\operatorname{Lef}(X,Y)$，则 $Y$ 的\Term{open-neighborhood}{开邻域} $U$ 的\Term{finite}{有限}\Term{flat}{平坦}\Term{covering}{覆盖} $R$ \Term{connected}{连通}，当且仅当 $R\times_U\widehat X$ \Term{connected}{连通}。特别地，$Y$ \Term{connected}{连通}，当且仅当 $Y$ 的\Term{open-neighborhood}{开邻域} $U$ \Term{connected}{连通}，也当且仅当 $X$ \Term{connected}{连通}。
\end{SGAresult}
\begin{proof}
\Term{local-ring}{局部环}\Term{ringed-space}{赋环空间} $(X,\mathcal O_X)$ \Term{connected}{连通}，当且仅当 $\Gamma(X,\mathcal O_X)$ 不能\Term{resolution}{分解}成两个非零环的直积。而由 $\operatorname{Lef}(X,Y)$ 有
\[\Gamma(U,r_*(\mathcal O_R))\simeq\Gamma(\widehat X,\widehat r_*(\mathcal O_{\widehat R})).\]
\end{proof}
\begin{SGAresult}{推论}{2.5}
若有 $\operatorname{Lef}(X,Y)$，则对每个 $U$，\Term{functor}{函子} $\operatorname{Et}(U)\to\operatorname{Et}(Y)$ \Term{fully-faithful}{全忠实}。若有 $\operatorname{Leff}(X,Y)$，则对 $Y$ 的每个\Term{etale}{平展}\Term{covering}{覆盖} $R$，存在 $Y$ 的\Term{open-neighborhood}{开邻域} $U$ 及 $U$ 的\Term{covering}{覆盖} $R'$，使 $R'\times_UY\simeq R$。
\end{SGAresult}
\begin{SGAresult}{推论}{2.6}
若有 $\operatorname{Lef}(X,Y)$ 且 $Y$ \Term{connected}{连通}，则 $Y$ 的每个\Term{open-neighborhood}{开邻域} $U$ \Term{connected}{连通}，自然\Term{homomorphism}{同态} $\pi_1(Y)\to\pi_1(U)$ \Term{surjectivity}{满射}。若此外有 $\operatorname{Leff}(X,Y)$，则自然\Term{homomorphism}{同态}
\[\pi_1(Y)\longrightarrow\varprojlim_U\pi_1(U)\]
是\Term{isomorphism}{同构}。（注意：为定义\Term{fundamental-group}{基本群}，假设已在 $Y$ 中选定一个“基点”，也以之作为 $X$ 的基点。）\footnote{编者注：结合 3.3 及判据 2.4、3.4，得到如下相对 Lefschetz 定理。设 $f:X\to S$ 是\Term{connected}{连通}\Term{noetherian}{诺特}\Term{scheme}{概形}之间的\Term{projective}{射影}\Term{flat}{平坦}\Term{morphism}{态射}，$D$ 是 $X$ 中\Term{effective}{有效}相对 Cartier \Term{divisor}{除子}且相对\Term{ample}{丰沛}。若对每个 $s\in S$，$X_s$ 在每个\Term{closed-point}{闭点}的\Term{depth}{深度}至少为 $2$，则 $D$ \Term{connected}{连通}，且对 $X$ 中包含 $D$ 的每个\Term{open-subset}{开集} $U$，\Term{arrow}{箭头} $i_U:\pi_1(D)\to\pi_1(U)$ \Term{surjectivity}{满射}。若此外 $X_s$ 沿 $D_s$ 每个\Term{closed-point}{闭点}的\Term{depth}{深度}至少为 $3$，且 $X$ 在\Term{closed-point}{闭点}处的\Term{local-ring}{局部环}是\Term{pure}{纯的}，例如为\Term{complete-intersection}{完全交}（参见 X 3.4），则 $i_X$ 是\Term{isomorphism}{同构}。参见 Bost J.-B., ``Lefschetz theorem for Arithmetic Surfaces'', Ann. Sci. Éc. Norm. Sup. (4) 32 (1999), pp.~241--312，定理 1.1、2.1。当 $X$ 仅为域上\Term{smooth}{光滑}几何\Term{connected}{连通}\Term{projective}{射影}曲面时，只要 $D$ 数值\Term{effective}{有效}且自交数正，仍有 $D$ \Term{connected}{连通}和 $\pi_1(D)\to\pi_1(U)$ \Term{surjectivity}{满射}（$U$ 是包含 $D$ 的\Term{open-subset}{开集}）；参见该文定理 2.3，及处理\Term{normal}{正规}\Term{complete}{完备}曲面的定理 2.4。当 $X$ 是整数环 $\mathcal O_K$ 上的\Term{normal}{正规}\Term{quasi-projective-morphism}{拟射影}算术曲面时，Bost 改进了 Ihara 的结果（Ihara Y., ``Horizontal divisors on arithmetic surfaces associated with Belyĭ uniformizations'', in The Grothendieck theory of dessins d'enfants (Luminy, 1993), London Math. Soc. Lect. Note Series, vol.~200, Cambridge Univ. Press, Cambridge, 1994, pp.~245--254，或前引文推论 7.2），证明：若点 $P\in X(\mathcal O_K)$（它在几何情形中扮演\Term{divisor}{除子} $D$ 的角色）满足某些正性条件，则由\Term{projection}{投影}得到的\Term{arrow}{箭头} $\pi_1(X)\to\pi_1(\operatorname{Spec}\mathcal O_K)$ \Term{invertible}{可逆}，其逆是由 $P$ 得到的\Term{arrow}{箭头} $\pi_1(\operatorname{Spec}\mathcal O_K)\to\pi_1(X)$（前引文定理 1.2）。}
\end{SGAresult}
这一切都直接由命题 1.1 和命题 2.3 得出。
\SourcePage{102}
\SGAsection{3}{比较基本群 $\pi_1(X)$ 与 $\pi_1(U)$}
\begin{SGAresult}{定义}{3.1}
设 $X$ 为\Term{prescheme}{预概形}，$Z$ 为 $X$ 的\Term{closed-subset}{闭子集}。置 $U=X-Z$。称对 $(X,Z)$ 是\Term{pure}{纯的}，如果对 $X$ 的每个\Term{open-subset}{开集} $V$，\Term{functor}{函子}
\[\operatorname{Et}(V)\longrightarrow\operatorname{Et}(V\cap U),\qquad V'\longmapsto V'\times_V(V\cap U)\]
是\Term{equivalence-of-categories}{范畴等价}。\footnote{关于在某些方面更令人满意的概念，参见评注 XIV 1.6 d)。}
\end{SGAresult}
\begin{SGAresult}{定义}{3.2}
设 $A$ 为\Term{noetherian}{诺特}\Term{local-ring}{局部环}，置 $X=\operatorname{Spec}A$。设 $\mathfrak r(A)$ 为 $A$ 的根，$x=\mathfrak r(A)$ 为 $X$ 的\Term{closed-point}{闭点}。称 $A$ 是\Term{pure}{纯的}，如果对 $(X,\{x\})$ 是\Term{pure}{纯的}。
\end{SGAresult}
我们把下面命题留给读者不去证明：
\begin{SGAresult}{命题}{3.3}
设 $X$ 为\Term{locally-noetherian}{局部诺特}\Term{prescheme}{预概形}，$Z$ 为 $X$ 的\Term{closed-subset}{闭子集}。对 $(X,Z)$ 是\Term{pure}{纯的}，当且仅当对每个 $z\in Z$，环 $\mathcal O_{X,z}$ 是\Term{pure}{纯的}。\footnote{比较 XIV 1.8 的非交换情形，其证明与 3.3 的证明本质相同。}
\end{SGAresult}
本节的主要结果是下列定理。
\begin{SGAresult}{定理}{3.4（\Term{purity-theorem}{纯性定理}）}
(i) \Term{dimension}{维数}至少为 $2$ 的\Term{regular}{正则}\Term{noetherian}{诺特}\Term{local-ring}{局部环}是\Term{pure}{纯的}（Zariski--Nagata \Term{purity}{纯性}定理）。
(ii) \Term{dimension}{维数}至少为 $3$ 且为\Term{complete-intersection}{完全交}的\Term{noetherian}{诺特}\Term{local-ring}{局部环}是\Term{pure}{纯的}。\footnote{编者注：关于所用方法的历史，见 Grothendieck 于 1961 年 10 月 1 日致 Serre 的信，Correspondance Grothendieck--Serre，Pierre Colmez、Jean-Pierre Serre 编，Documents Mathématiques, vol.~2, Société Mathématique de France, Paris, 2001。}
\end{SGAresult}
\SourcePage{103}
回忆：称\Term{local-ring}{局部环}为\Term{complete-intersection}{完全交}，是指存在\Term{regular}{正则}\Term{noetherian}{诺特}\Term{local-ring}{局部环} $B$，以及由 $B$ 的根 $\mathfrak r(B)$ 中元素组成的 $B$-\Term{regular}{正则}序列 $(t_1,\ldots,t_k)$，使
\[A\simeq B/(t_1,\ldots,t_k).\]
这里说 $A$ 是绝对\Term{complete-intersection}{完全交}歧义会更小，以区别于先前遇到的情形：$X$ 为\Term{locally-noetherian}{局部诺特}\Term{prescheme}{预概形}（未必\Term{regular}{正则}），$Y$ 是 $X$ 的\Term{closed-subset}{闭子集}，而我们说 $Y$“在 $X$ 中\Term{local-scheme}{局部}在\Term{set}{集合}意义下是\Term{complete-intersection}{完全交}”。先证明几个引理。
\begin{SGAresult}{引理}{3.5}
设 $X$ 为\Term{locally-noetherian}{局部诺特}\Term{prescheme}{预概形}，$U$ 为其\Term{open-subset}{开子集}，置 $Z=X-U$。设 $i:U\to X$ 为\Term{canonical}{典范}\Term{immersion}{浸入}。下列条件\Term{equivalence}{等价}：
(i) 对 $X$ 的每个\Term{open-subset}{开集} $V$，置 $V'=V\cap U$，则从\Term{locally-free}{局部自由}\Term{coherent}{凝聚} $\mathcal O_V$-模\Term{category}{范畴}到\Term{locally-free}{局部自由}\Term{coherent}{凝聚} $\mathcal O_{V'}$-模\Term{category}{范畴}的\Term{functor}{函子} $F\mapsto F|_{V'}$ \Term{fully-faithful}{全忠实}；
(ii) 自然\Term{homomorphism}{同态} $\mathcal O_X\to i_*(\mathcal O_U)$ 是\Term{isomorphism}{同构}；
(iii) 对每个 $z\in Z$，有 $\operatorname{prof}\mathcal O_{X,z}\geq2$。
\end{SGAresult}
\begin{proof}
(ii)、(iii) 的\Term{equivalence}{等价}性已见于 III 3.3。证明 (ii) 蕴涵 (i)：若 $F,G$ 是两个\Term{locally-free}{局部自由}\Term{coherent}{凝聚} $\mathcal O_V$-模，$\mathcal Hom(F,G)$ 也如此，故 $\mathcal Hom(F,G)\to i_*(\mathcal Hom(F|_{V'},G|_{V'}))$ 为\Term{isomorphism}{同构}，所以 $\operatorname{Hom}(F,G)\simeq\operatorname{Hom}(F|_{V'},G|_{V'})$。反过来，取 $F=G=\mathcal O_X$，对 $X$ 的每个\Term{open-subset}{开集} $V$ 应用 (i)。
\end{proof}
下面是一个有用的“\Term{descent}{下降}引理”。
\begin{SGAresult}{引理}{3.6}
设 $X$ 为\Term{locally-noetherian}{局部诺特}\Term{prescheme}{预概形}，$Z$ 为 $X$ 的\Term{closed-subset}{闭子集}，置 $U=X-Z$。假设\Term{homomorphism}{同态} $\mathcal O_X\to i_*(\mathcal O_U)$ 是\Term{isomorphism}{同构}。设 $f:X_1\to X$ 为\Term{faithfully-flat}{忠实平坦}\Term{quasi-compact}{拟紧}\Term{morphism}{态射}，置 $Z_1=f^{-1}(Z)$。若对 $(X_1,Z_1)$ 是\Term{pure}{纯的}，则 $(X,Z)$ 也是\Term{pure}{纯的}。
\end{SGAresult}
\begin{proof}
注意假设 $\mathcal O_X\simeq i_*(\mathcal O_U)$ 在\Term{flat}{平坦}\Term{base-extension}{基扩张}下保持，因为 $i$ 为\Term{quasi-compact}{拟紧}\Term{morphism}{态射}，这时\Term{direct-image}{直像}与\Term{inverse-image}{逆像}可交换。将\Term{etale}{平展}\Term{covering}{覆盖}解释为\Term{locally-free}{局部自由}\Term{coherent}{凝聚}\Term{algebra}{代数}，可知此假设意味着\Term{functor}{函子}
\[\operatorname{Et}(V)\longrightarrow\operatorname{Et}(U\cap V),\qquad V'\longmapsto V'\times_V(V\cap U)\]
\Term{fully-faithful}{全忠实}。只须证明\Term{effective}{有效}性。例如可引入 $X_1$ 在 $X$ 上的平方 $X_2$ 和立方 $X_3$，并注意\Term{faithfully-flat}{忠实平坦}\Term{quasi-compact}{拟紧}\Term{morphism}{态射}对于\Term{prescheme}{预概形}\Term{category}{范畴}上的\Term{etale}{平展}\Term{covering}{覆盖}\Term{fibre-category}{纤维范畴}是\Term{universal}{普遍}\Term{effective-descent}{有效下降}\Term{morphism}{态射}。
\SourcePage{104}
从这里形式地得到结论。\footnote{参见 J. Giraud, Méthode de la descente, Mémoire no.~2 du Bulletin de la Société Mathématique de France (1964)。}
\end{proof}
\textbf{注记 3.7。} 在此过程中我们已证明：若 $\mathcal O_X\to i_*(\mathcal O_U)$ 是\Term{isomorphism}{同构}，则 $X$ \Term{connected}{连通}当且仅当 $U$ \Term{connected}{连通}，并且此时 $\pi_1(U)\to\pi_1(X)$ \Term{surjectivity}{满射}。
\begin{SGAresult}{推论}{3.8}
设 $A$ 为\Term{noetherian}{诺特}\Term{local-ring}{局部环}，假设 $\operatorname{prof}A\geq2$。若 $\widehat A$ 是\Term{pure}{纯的}，则 $A$ 是\Term{pure}{纯的}。
\end{SGAresult}
由引理 3.5、3.6 得出。下面的引理是\Term{purity}{纯性}定理证明的关键。
\begin{SGAresult}{引理}{3.9}
设 $A$ 为\Term{noetherian}{诺特}\Term{local-ring}{局部环}，$t\in\mathfrak r(A)$ 为 $A$-\Term{regular}{正则}元。假设 $A$ 对 $t$-\Term{adic-topology}{进拓扑}\Term{complete}{完备}，且为\Term{regular}{正则}\Term{local-ring}{局部环}的商（例如 $A$ \Term{complete}{完备}）。置 $B=A/tA$。
(i) 若对 $A$ 的每个满足 $\dim A/\mathfrak p=1$ 的\Term{prime-ideal}{素理想} $\mathfrak p$，有 $\operatorname{prof}A_{\mathfrak p}\geq2$，则 $B$ 纯蕴涵 $A$ 纯。
(ii) 若对 $A$ 的每个满足 $\dim A/\mathfrak p=1$ 的\Term{prime-ideal}{素理想} $\mathfrak p$，有 $\operatorname{prof}A_{\mathfrak p}\geq2$，当 $t\notin\mathfrak p$ 时 $A_{\mathfrak p}$ 纯，且当 $t\in\mathfrak p$ 时 $\operatorname{prof}A_{\mathfrak p}\geq3$，则 $A$ 纯蕴涵 $B$ 纯。\footnote{编者注：最后一个条件可以改进，参见原书第 90 页编者注 (1)。}
\end{SGAresult}
\begin{proof}
设 $X'=\operatorname{Spec}(A)$，$Y'=V(t)$，并将后者与 $B$ 的谱等同。设 $x=\mathfrak r(A)$，置 $X=X'-\{x\}$、$Y=Y'-\{x\}=X\cap Y'$。以 $\widehat{X'}$ 表示 $A$ 对 $t$-\Term{adic-topology}{进拓扑}的形式谱，它可等同于 $X'$ 沿 $Y'$ 的\Term{formal-completion}{形式完备化}。因 $A$ 对 $t$-\Term{adic-topology}{进拓扑}\Term{complete}{完备}，$\operatorname{Et}(X')\to\operatorname{Et}(\widehat{X'})$ 是\Term{equivalence-of-categories}{范畴等价}。又由命题 1.1，$\operatorname{Et}(\widehat{X'})\to\operatorname{Et}(Y')$ 是\Term{equivalence}{等价}，所以 $\operatorname{Et}(X')\to\operatorname{Et}(Y')$ 是\Term{equivalence-of-categories}{范畴等价}。
证明 (i)。考虑图
\[\begin{tikzcd}X'&X\arrow[l]\\Y'\arrow[u]&Y\arrow[l]\arrow[u]\end{tikzcd}\qquad
\begin{tikzcd}\operatorname{Et}(X')\arrow[r,"a"]\arrow[d,"c"']&\operatorname{Et}(X)\arrow[d,"b"]\\\operatorname{Et}(Y')\arrow[r,"d"]&\operatorname{Et}(Y).\end{tikzcd}\]
刚已证明 $c$ 是\Term{equivalence}{等价}；由 $B$ \Term{pure}{纯的}假设，$d$ 也是\Term{equivalence}{等价}；最后由例 2.1（参见 2.3 (i)），$b$ \Term{fully-faithful}{全忠实}。
证明 (ii)。这次假设 $A$ 纯，所以 $a$ 为\Term{equivalence}{等价}，$c$ 也如此。证明 $b$ 为\Term{equivalence}{等价}。由例 2.1 已知 $\operatorname{Leff}(X,Y)$，所以 $b$ 已\Term{fully-faithful}{全忠实}，只须证明它\Term{essentially-surjective}{本质满}。应用 2.3 (ii)，并注意若 $U$ 为 $Y$ 在 $X$ 中的\Term{open-neighborhood}{开邻域}，
\SourcePage{105}
则 $U$ 在 $X$ 中的补集是\Term{finite}{有限}个\Term{closed-point}{闭点}的并；由命题 3.3，对 $(X,X-U)$ 是\Term{pure}{纯的}，因为在这样的点 $p$，按假设 $\mathcal O_{X,p}=A_{\mathfrak p}$ 是\Term{pure}{纯的}。结论由此得到。
\end{proof}
\begin{proof}[\Term{purity}{纯性}定理的证明]
先对\Term{dimension}{维数}归纳证明 (i)。设 $A$ 为\Term{dimension}{维数} $2$ 的\Term{noetherian}{诺特}\Term{local-ring}{局部环}。置 $X'=\operatorname{Spec}(A)$、$x=\mathfrak r(A)$、$X=X'-\{x\}$。有 $\operatorname{prof}A=2$，故可对 $(X',\{x\})$ 应用引理 3.5，从而 $\operatorname{Et}(X')\to\operatorname{Et}(X)$ \Term{fully-faithful}{全忠实}。现在设 $r:R\to X$ 为\Term{etale}{平展}\Term{covering}{覆盖}，由\Term{locally-free}{局部自由}\Term{coherent}{凝聚}\Term{etale}{平展} $\mathcal O_X$-\Term{algebra}{代数} $\mathcal A=r_*(\mathcal O_R)$ 定义。以 $i:X\to X'$ 表示\Term{canonical}{典范}\Term{immersion}{浸入}。我断言 $i_*(\mathcal A)=\mathcal B$ 是\Term{coherent}{凝聚} $\mathcal O_{X'}$-\Term{algebra}{代数}；事实上只须应用“\Term{finiteness-theorem}{有限性定理}” VIII 2.3。又断言此\Term{algebra}{代数}在 $x$ 的\Term{depth}{深度}至少为 $2$：这是因为它是 $\mathcal O_X$-模的\Term{direct-image}{直像}，且 $X=X'-\{x\}$。因为 $A$ 是\Term{dimension}{维数} $2$ 的\Term{regular}{正则}环，
\[\operatorname{dp}\mathcal B+\operatorname{prof}\mathcal B=\dim A=2,\]
其中 $\operatorname{dp}\mathcal B$ 表示 $\mathcal B$ 的\Term{projective-dimension}{投射维数}。于是 $\operatorname{dp}\mathcal B=0$，故 $\mathcal B$ \Term{projective-module}{投射}，从而自由。所以 $\mathcal B$ 定义 $X'=\operatorname{Spec}(A)$ 的\Term{finite}{有限}\Term{flat}{平坦}\Term{covering}{覆盖}。此\Term{covering}{覆盖}不\Term{etale}{平展}的点集是 $X'$ 的\Term{closed-subset}{闭子集}，由主\Term{ideal}{理想}，即 $\mathcal B/A$ 的\Term{discriminant}{判别式}\Term{ideal}{理想}定义。按构造，此闭集包含在 $x=\mathfrak r(A)$ 中，因 $\dim A=2$，它必为空。
设 $A$ 为\Term{regular}{正则}\Term{noetherian}{诺特}\Term{local-ring}{局部环}，$\dim A=n\geq3$。假设 (i) 对\Term{dimension}{维数}小于 $n$ 的环已证。为证明 $A$ 纯，由 3.8 可假设 $A$ \Term{complete}{完备}。取 $t\in\mathfrak r(A)$，使其在 $\mathfrak r(A)/\mathfrak r(A)^2$ 中的像非零。则 $B=A/tA$ 是\Term{dimension}{维数} $n-1$ 的\Term{regular}{正则}\Term{noetherian}{诺特}\Term{local-ring}{局部环}，因为 $n-1\geq2$，所以它纯。因 $A$ \Term{complete}{完备}，可应用引理 3.9 (i)，即得结论。
证明 (ii)。设 $A$ 是\Term{dimension}{维数}至少为 $3$ 的\Term{noetherian}{诺特}\Term{local-ring}{局部环}。假设存在\Term{regular}{正则}\Term{noetherian}{诺特}\Term{local-ring}{局部环} $B$ 及 $B$-序列 $(t_1,\ldots,t_k)$，使 $A\simeq B/(t_1,\ldots,t_k)$。对 $k$ 归纳证明 $A$ 纯。若 $k=0$，由 (i) 已知。设 $k\geq1$，且结果对 $k'<k$ 已知。由推论 3.9 可假设 $A$（因而 $B$ 也）\Term{complete}{完备}。置 $C=B/(t_1,\ldots,t_{k-1})$，则 $A\simeq C/t_kC$，$t_k$ 是 $C$-\Term{regular}{正则}元。归纳假设给出 $C$ 纯，只须证明可应用引理 3.9 (ii)。引理中的 $A,B$ 分别变成 $C,A$。有 $\dim C\geq4$，所以对 $C$ 中满足 $\dim C/\mathfrak p=1$ 的每个\Term{prime-ideal}{素理想} $\mathfrak p$，有 $\operatorname{prof}C_{\mathfrak p}\geq3$。此外 $C_{\mathfrak p}$ 是 $k'\leq k-1$ 的\Term{complete-intersection}{完全交}，故由归纳假设为纯。
\end{proof}
\begin{SGAresult}{定理}{3.10}
设 $X$ 为\Term{locally-noetherian}{局部诺特}\Term{prescheme}{预概形}，$Y$ 为 $X$ 的\Term{closed-subset}{闭子集}。假设有 $\operatorname{Leff}(X,Y)$（参见例 2.1、2.2）。再假设：对 $Y$ 的每个\Term{open-neighborhood}{开邻域} $U$ 及每个 $x\in X-U$，\Term{local-ring}{局部环} $\mathcal O_{X,x}$ 或者是\Term{dimension}{维数}至少为 $2$ 的\Term{regular}{正则}环，或者是\Term{dimension}{维数}至少为 $3$ 的\Term{complete-intersection}{完全交}。则
\[\pi_0(Y)\longrightarrow\pi_0(X)\]
\SourcePage{106}
是\Term{bijection}{双射}；若 $X$ \Term{connected}{连通}，则 $\pi_1(Y)\to\pi_1(X)$ 是\Term{isomorphism}{同构}。
\end{SGAresult}
已无须证明。注意在所引的例 2.1、2.2 中，$U$ 的补集都是\Term{finite}{有限}个\Term{closed-point}{闭点}的并，因此关于 $\mathcal O_{X,x}$ \Term{dimension}{维数}的假设并非无稽之谈。
\SourcePage{107}
\SGAchapter{XI}{对 Picard 群的应用}{Application au groupe de Picard}
本讲仿照上一讲，但这次第 1 节的结果较弱。本讲始终以 $X$ 表示\Term{locally-noetherian}{局部诺特}\Term{prescheme}{预概形}，$\mathcal I$ 表示 $\mathcal O_X$ 的\Term{quasi-coherent}{拟凝聚}\Term{ideal}{理想}（因而 $Y=V(\mathcal I)$ 是 $X$ 的\Term{closed-subset}{闭子集}），$U$ 表示 $Y$ 在 $X$ 中的可变\Term{open-neighborhood}{开邻域}，$\widehat X$ 表示 $X$ 沿 $Y$ 的\Term{formal-completion}{形式完备化}。对任意\Term{ringed-space}{赋环空间} $(Z,\mathcal O_Z)$，以 $P(Z)$ 表示\Term{invertible}{可逆} $\mathcal O_Z$-模，即秩 $1$ \Term{locally-free}{局部自由}模的\Term{category}{范畴}，以 $\operatorname{Pic}(Z)$ 表示 $Z$ 上\Term{invertible}{可逆}模的\Term{isomorphism}{同构}类群。
\SGAsection{1}{比较完备化的 Picard 群与 Pic(Y)}
对每个 $n\in\mathbb N$，置 $X_n=(Y,\mathcal O_X/\mathcal I^{n+1})$、$P_n=\mathcal I^{n+1}/\mathcal I^{n+2}$。$Y$ 上的\Term{abelian-group}{阿贝尔群}\Term{sheaf}{层}序列
\[0\longrightarrow P_n\xrightarrow{u}\mathcal O^*_{X_{n+1}}\xrightarrow{v}\mathcal O^*_{X_n}\longrightarrow1\tag{1.1}\]
\Term{exact}{正合}。具体说，$P_n$ 的群结构为加法结构，对每个 $x\in P_n$ 有 $u(x)=1+x$，$v$ 是由包含 $\mathcal I^{n+2}\to\mathcal I^{n+1}$ 得到的\Term{homomorphism}{同态}。对每个 $y\in Y$，$\mathcal O_{X_n,y}$ 是\Term{local-ring}{局部环} $\mathcal O_{X_{n+1},y}$ 模掉\Term{nilpotent}{幂零}\Term{ideal}{理想}的商，由此可见 $v$ \Term{surjectivity}{满射}；其余同样显然。从 (1.1) 得\Term{cohomology}{上同调}\Term{exact-sequence}{正合列}
\[H^1(Y,P_n)\xrightarrow{u^1}H^1(Y,\mathcal O^*_{X_{n+1}})\xrightarrow{v^1}H^1(Y,\mathcal O^*_{X_n})\xrightarrow{d}H^2(Y,P_n).\tag{*}\]
另一方面，对每个 $n\in\mathbb N$，已知可将 $\operatorname{Pic}(X_n)$ 与 $H^1(Y,\mathcal O^*_{X_n})$ 等同；此外，若\Term{invertible}{可逆} $\mathcal O_{X_{n+1}}$-模 $E$ 对应于\Term{cohomology}{上同调}类 $c(E)$，则 $E$ 在 $X_n$ 上的\Term{inverse-image}{逆像}对应的\Term{cohomology}{上同调}类等于 $v^1(c(E))$。得到如下命题。
\SourcePage{108}
\begin{SGAresult}{命题}{1.1}
保持上述记号。设 $p\in\mathbb N$。映射 $\operatorname{Pic}(\widehat X)\to\operatorname{Pic}(Y_n)$：
(i) 若对 $n\geq p$ 有 $H^1(Y,P_n)=0$，则在 $n\geq p$ 时为\Term{injectivity}{单射}；
(ii) 若对 $n\geq p$ 及 $i=1,2$ 有 $H^i(Y,P_n)=0$，则在 $n\geq p$ 时为\Term{isomorphism}{同构}。
\end{SGAresult}
当然，\Term{exact-sequence}{正合列} (*) 包含比上述命题更多的信息。读者会注意到，我们还没有讨论\Term{functor}{函子} $P(\widehat X)\to P(Y)$。给定两个\Term{invertible}{可逆} $\mathcal O_{\widehat X}$-模 $E,F$，$H=\mathcal Hom(E,F)$ 也\Term{invertible}{可逆}。若以下标 $n$ 表示模 $\mathcal I^{n+1}$ 的\Term{reduced}{约化}，则有\Term{exact-sequence}{正合列}
\[0\longrightarrow H_0\otimes P_n\longrightarrow\mathcal Hom(E_{n+1},F_{n+1})\longrightarrow\mathcal Hom(E_n,F_n)\longrightarrow0.\]
由此得一个\Term{cohomology}{上同调}\Term{exact-sequence}{正合列}，我们不再写出，其解释显而易见；可用这一观察研究\Term{functor}{函子} $P$。
\SGAsection{2}{比较 Pic(X) 与完备化的 Picard 群}
下列结果的证明可见第 X 讲第 2 节。
\begin{SGAresult}{命题}{2.1}
假设有 $\operatorname{Lef}(X,Y)$；则对 $Y$ 在 $X$ 中的每个\Term{open-neighborhood}{开邻域} $U$，\Term{functor}{函子}
\[P(U)\longrightarrow P(\widehat X)\tag{2.1}\]
\Term{fully-faithful}{全忠实}，所以映射
\[\operatorname{Pic}(U)\longrightarrow\operatorname{Pic}(\widehat X)\tag{2.2}\]
为\Term{injectivity}{单射}。若有 $\operatorname{Leff}(X,Y)$，则映射 (2.3) 为\Term{isomorphism}{同构}：
\[\varinjlim_U\operatorname{Pic}(U)\longrightarrow\operatorname{Pic}(\widehat X).\tag{2.3}\]
\end{SGAresult}
\begin{SGAresult}{推论}{2.2}
假设有 $\operatorname{Lef}(X,Y)$，且对每个整数 $n\geq p$ 有 $H^1(Y,P_n)=0$。则对每个\Term{open-subset}{开集} $U\supset Y$，映射
\[\operatorname{Pic}(X)\longrightarrow\operatorname{Pic}(U)\longrightarrow\operatorname{Pic}(Y_n)\]
在 $n\geq p$ 时为\Term{injectivity}{单射}。若有 $\operatorname{Leff}(X,Y)$，且此外对每个整数 $n\geq p$ 及 $i=1,2$ 有 $H^i(Y,P_n)=0$，则映射
\[\varinjlim_U\operatorname{Pic}(U)\longrightarrow\operatorname{Pic}(Y_n)\]
在 $n\geq p$ 时为\Term{isomorphism}{同构}。
\end{SGAresult}
\SourcePage{109}
\SGAsection{3}{比较 P(X) 与 P(U)}
先给一个定义。
\begin{SGAresult}{定义}{3.1}
设 $X$ 为\Term{prescheme}{预概形}，$Z$ 为 $X$ 的\Term{closed-subset}{闭子集}，置 $U=X-Z$。称 $X$ 在 $Z$ 的点处是\Term{parafactorial}{旁析因}的，如果对 $X$ 的每个\Term{open-subset}{开集} $V$，\Term{functor}{函子} $P(V)\to P(V\cap U)$ 是\Term{equivalence-of-categories}{范畴等价}。也称对 $(X,Z)$ 是\Term{parafactorial}{旁析因}的。\footnote{关于\Term{parafactorial}{旁析因}性更详细的研究及 3.3 的证明，参见 EGA IV 21.13、21.14。}
\end{SGAresult}
回忆 $P(Z)$ 表示 $Z$ 上秩 $1$ \Term{locally-free}{局部自由}模的\Term{category}{范畴}。
\begin{SGAresult}{定义}{3.2}
称\Term{noetherian}{诺特}\Term{local-ring}{局部环}是\Term{parafactorial}{旁析因}的，如果对 $(\operatorname{Spec}(A),\{\mathfrak r(A)\})$ 是\Term{parafactorial}{旁析因}的。
\end{SGAresult}
下面的命题说明该概念是“逐点的”。
\begin{SGAresult}{命题}{3.3}
假设 $X$ \Term{locally-noetherian}{局部诺特}。对 $(X,Z)$ \Term{parafactorial}{旁析因}，当且仅当对每个 $z\in Z$，\Term{local-ring}{局部环} $\mathcal O_{X,z}$ \Term{parafactorial}{旁析因}。
\end{SGAresult}
注意\Term{parafactorial}{旁析因}蕴含“\Term{fully-faithful}{全忠实}”。仿照第 X 讲引理 3.5 可证：
\begin{SGAresult}{引理}{3.4}
若 $X$ 为\Term{locally-noetherian}{局部诺特}\Term{prescheme}{预概形}，$Z=X-U$ 为 $X$ 的\Term{closed-subset}{闭子集}，则下列条件\Term{equivalence}{等价}：
(i) 对 $X$ 的每个\Term{open-subset}{开集} $V$，\Term{functor}{函子} $P(V)\to P(V\cap U)$ \Term{fully-faithful}{全忠实}；
(ii) \Term{homomorphism}{同态} $\mathcal O_X\to i_*(\mathcal O_U)$ 是\Term{isomorphism}{同构}；
(iii) 对每个 $z\in Z$，有 $\operatorname{prof}(\mathcal O_{X,z})\geq2$。
\end{SGAresult}
所以“\Term{parafactorial}{旁析因}”意味着 3.4 的条件成立，并且对 $X$ 的每个\Term{open-subset}{开集} $V$，\Term{homomorphism}{同态} $\operatorname{Pic}(V)\to\operatorname{Pic}(V\cap U)$ \Term{surjectivity}{满射}。特别地，当 $X$ 是\Term{noetherian}{诺特}\Term{local-ring}{局部环}的谱时得到：
\begin{SGAresult}{命题}{3.5}
设 $A$ 为\Term{noetherian}{诺特}\Term{local-ring}{局部环}。$A$ \Term{parafactorial}{旁析因}的充要条件是 $\operatorname{prof}A\geq2$ 且 $\operatorname{Pic}(X'-\{x\})=0$，其中 $X'=\operatorname{Spec}(A)$，$x$ 为 $X'$ 的唯一\Term{closed-point}{闭点}。
\end{SGAresult}
注意\Term{dimension}{维数}不超过 $1$ 的\Term{local-ring}{局部环}从来不是\Term{parafactorial}{旁析因}的，因为其\Term{depth}{深度}不超过 $1$。所以\Term{factorial}{析因}不蕴含\Term{parafactorial}{旁析因}；不过，对\Term{dimension}{维数}至少为 $2$ 的\Term{noetherian}{诺特}\Term{local-ring}{局部环}，这个蕴涵成立，下文即将说明。
\SourcePage{110}
\begin{SGAresult}{引理}{3.6}
设 $X$ 为\Term{locally-noetherian}{局部诺特}\Term{prescheme}{预概形}，$Z$ 为 $X$ 的\Term{closed-subset}{闭子集}。设 $f:X_1\to X$ 为\Term{faithfully-flat}{忠实平坦}\Term{quasi-compact}{拟紧}\Term{morphism}{态射}，置 $Z_1=f^{-1}(Z)$。若 $(X_1,Z_1)$ \Term{parafactorial}{旁析因}，则 $(X,Z)$ \Term{parafactorial}{旁析因}。
\end{SGAresult}
\begin{proof}
先注意，若 $i:(X-Z)\to X$ 是 $U=X-Z$ 到 $X$ 的\Term{canonical}{典范}\Term{immersion}{浸入}，则\Term{quasi-coherent}{拟凝聚} $\mathcal O_U$-模沿 $i$ 的\Term{direct-image}{直像}的构造与\Term{base-change}{基变换} $f$ 可交换，因为 $f$ \Term{flat}{平坦}。因此，对 $(X,Z)$ 或对 $(X_1,Z_1)$ 假设引理 3.5 的\Term{equivalence}{等价}条件是一样的，因为 $f$ 对\Term{quasi-coherent}{拟凝聚}\Term{sheaf}{层}\Term{category}{范畴}是\Term{descent}{下降}\Term{morphism}{态射}。只须证明：对 $X$ 的每个\Term{open-subset}{开集} $V$，$\operatorname{Pic}(V)\to\operatorname{Pic}(V\cap U)$ \Term{surjectivity}{满射}。作不改变任何事情的\Term{base-change}{基变换} $V\to X$（原文如此），便归结到 $V=X$。若 $L$ 是\Term{invertible}{可逆} $\mathcal O_U$-模，且有\Term{locally-free}{局部自由}延拓，则由刚才的论证，此延拓\Term{isomorphism}{同构}于 $i_*(L)$。只须证明 $i_*(L)$ \Term{invertible}{可逆}。再次利用沿 $i$ 的\Term{direct-image}{直像}与\Term{flat}{平坦}\Term{base-change}{基变换}可交换，以及“秩 $1$ \Term{locally-free}{局部自由}”可沿\Term{faithfully-flat}{忠实平坦}\Term{quasi-compact}{拟紧}\Term{morphism}{态射}\Term{descent}{下降}，便告完成。
\end{proof}
\begin{SGAresult}{推论}{3.7}
设 $A$ 为\Term{noetherian}{诺特}\Term{local-ring}{局部环}。若 $\widehat A$ \Term{parafactorial}{旁析因}，则 $A$ \Term{parafactorial}{旁析因}。
\end{SGAresult}
不要以为 $A$ \Term{parafactorial}{旁析因}就意味着 $\widehat A$ \Term{parafactorial}{旁析因}。\footnote{编者注：关于 $A$ 及其\Term{completion}{完备化}的\Term{factorial}{析因}性之间的联系的精确研究，见 Heitmann R., ``Characterization of completions of unique factorization domains'', Trans. Amer. Math. Soc. 337 (1993), no.~1, pp.~379--387。}
在陈述本节主要定理之前，先联系\Term{divisor}{除子}理论及\Term{factorial}{析因}环的概念。\footnote{关于下面的一般事实，另见 EGA IV 21。} 设 $X$ 为\Term{noetherian}{诺特}\Term{normal}{正规}\Term{prescheme}{预概形}。以 $Z^1(X)$ 表示由满足 $\dim\mathcal O_{X,x}=1$ 的 $x\in X$ 生成的自由\Term{abelian-group}{阿贝尔群}。这样的点的\Term{local-ring}{局部环}是\Term{discrete-valuation-ring}{离散赋值环}，以 $v_x$ 表示对应的规范化赋值。设 $K(X)$ 是 $X$ 的\Term{rational-function-ring}{有理函数环}，设
\[p:K(X)^*\longrightarrow Z^1(X)\]
把每个 $f\in K(X)^*$ 对应到余维 $1$ 的\Term{cycle}{闭链}
\[(f)=\sum_{x\in X,\ \dim\mathcal O_{X,x}=1}v_x(f)\cdot x.\]
$p$ 的像记作 $P(X)$，其元素称为\Term{principal-divisor}{主除子}。\footnote{按照 EGA IV 21 的术语，我们现在更愿意只把“\Term{divisor}{除子}”一词用于“\Term{locally-principal}{局部主}\Term{divisor}{除子}”或“Cartier \Term{divisor}{除子}”。} 置
\[\operatorname{Cl}(X)=Z^1(X)/P(X).\]
\SourcePage{111}
设 $Z'^1(X)$ 是 $Z^1(X)$ 中由\Term{locally-principal}{局部主}\Term{divisor}{除子}组成的\Term{subgroup}{子群}。已知
\[\operatorname{Pic}(X)\simeq Z'^1(X)/P(X),\]
所以 $\operatorname{Pic}(X)$ 可等同于 $\operatorname{Cl}(X)$ 的\Term{subgroup}{子群}。
注意若 $U$ 是 $X$ 的\Term{dense}{稠密}\Term{open-subset}{开集}，$K(X)\to K(U)$ 是\Term{isomorphism}{同构}；若 $\operatorname{codim}(X-U,X)\geq2$，即每个满足 $\dim\mathcal O_{X,x}\leq1$ 的 $x\in X$ 都属于 $U$，则\Term{homomorphism}{同态} $Z^1(X)\to Z^1(U)$、从而 $\operatorname{Cl}(X)\to\operatorname{Cl}(U)$ 也是\Term{isomorphism}{同构}。最后，若每个 $x\in U$ 都\Term{factorial}{析因}，即 $\mathcal O_{X,x}$ \Term{factorial}{析因}，则 $Z^1(U)=Z'^1(U)$，故 $\operatorname{Pic}(U)\simeq\operatorname{Cl}(U)$。
\begin{SGAresult}{命题}{3.7.1}
设 $X$ 为\Term{noetherian}{诺特}\Term{normal}{正规}\Term{prescheme}{预概形}，$(U_i)_{i\in I}$ 为 $X$ 的一族\Term{open-subset}{开集}，满足：
a) $U_i$ 构成滤子基；\footnote{编者注：即递减\Term{filtration}{滤过}族。}
b) 置 $Y_i=X-U_i$，则对每个 $i$ 有 $\operatorname{codim}(Y_i,X)\geq2$；
c) 若 $x\in U_i$ 对所有 $i\in I$ 成立，则 $\mathcal O_{X,x}$ \Term{factorial}{析因}。
则有\Term{isomorphism}{同构}
\[\varinjlim_{i\in I}\operatorname{Pic}(U_i)\xrightarrow\sim\operatorname{Cl}(X).\]
\end{SGAresult}
\begin{proof}
b) 意味着满足 $\dim\mathcal O_{X,x}\leq1$ 的每个 $x\in X$ 都属于所有 $U_i$。因此 $U_i$ \Term{dense}{稠密}，且\Term{homomorphism}{同态} $Z^1(U_i)\to Z^1(X)$ 与 $K(U_i)\to K(X)$ 都是\Term{isomorphism}{同构}。所以 $\operatorname{Pic}(U)\subset\operatorname{Cl}(U_i)\simeq\operatorname{Cl}(X)$。只须证明每个 $D\in Z^1(X)$ 都对某个适当的 $i$ 属于 $Z'^1(U_i)$。只须对\Term{irreducible}{不可约}\Term{effective}{有效}“\Term{divisor}{除子}”证明。于是设 $x\in X$，$\dim\mathcal O_{X,x}=1$。只须证明存在 $i\in I$，使 $\overline{\{x\}}$ 在 $U_i$ 的各点\Term{locally-principal}{局部主}。设 $\mathcal I$ 是闭集 $\overline{\{x\}}$ 的最大\Term{ideal-of-definition}{定义理想}。$\mathcal I$ 在其邻域上自由的点组成\Term{open-subset}{开集} $U$。由 c)，$U\supset\bigcap_{i\in I}U_i$。若置 $Y=X-U$，则 $Y\subset\bigcup_{i\in I}Y_i$，其中 $Y_i=X-U_i$。因 $Y$ 闭，它只有\Term{finite}{有限}个\Term{generic-point}{一般点}，所以包含在\Term{finite}{有限}个 $Y_i$ 的并中；由于 $U_i$ 构成滤子基，便包含在某个 $Y_j$ 中。因此 $U\supset U_j$。
\end{proof}
\begin{SGAresult}{推论}{3.8}
设 $X$ 为\Term{noetherian}{诺特}\Term{normal}{正规}\Term{prescheme}{预概形}，$Y$ 为余维至少为 $2$ 的\Term{closed-subset}{闭子集}。假设对每个 $p\in X-Y$，$\mathcal O_{X,p}$ \Term{factorial}{析因}，则
\[\operatorname{Pic}(X-Y)\longrightarrow\operatorname{Cl}(X-Y)\longrightarrow\operatorname{Cl}(X)\]
均为\Term{isomorphism}{同构}。
\end{SGAresult}
\SourcePage{112}
\begin{SGAresult}{推论}{3.9}
设 $A$ 为\Term{noetherian}{诺特}\Term{normal}{正规}\Term{local-ring}{局部环}，置 $X'=\operatorname{Spec}(A)$、$x=\mathfrak r(A)$。$A$ \Term{factorial}{析因}，当且仅当 $\operatorname{Pic}(X'-\{x\})=0$ 且对每个 $\mathfrak p\in X'-\{x\}$，$A_{\mathfrak p}$ \Term{factorial}{析因}。
\end{SGAresult}
事实上，$A$ \Term{factorial}{析因}的充要条件是 $\operatorname{Cl}(X')=0$。\footnote{编者注：见 Bourbaki, Algèbre commutative VII.1.4，定理 2 的推论，以及 VII.3.2，定理 1。}
\begin{SGAresult}{推论}{3.10}
设 $A$ 为\Term{dimension}{维数}至少为 $2$ 的\Term{noetherian}{诺特}\Term{local-ring}{局部环}。置 $X'=\operatorname{Spec}(A)$、$x=\mathfrak r(A)$、$X=X'-\{x\}$。下列条件\Term{equivalence}{等价}：
(i) $A$ \Term{factorial}{析因}；
(ii) a) 对每个 $y\in X$，$\mathcal O_{X,p}$ \Term{factorial}{析因}；b) $A$ \Term{parafactorial}{旁析因}，即 $\operatorname{prof}A\geq2$ 且 $\operatorname{Pic}(X)=0$。
\end{SGAresult}
证明此推论之前，先陈述：
\begin{SGAresult}{Serre \Term{normality}{正规性}判据}{3.11}
设 $A$ 为\Term{noetherian}{诺特}\Term{local-ring}{局部环}。$A$ \Term{normal}{正规}的充要条件是：
(i) 对 $A$ 的每个满足 $\dim A_{\mathfrak p}\leq1$ 的\Term{prime-ideal}{素理想} $\mathfrak p$，$A_{\mathfrak p}$ \Term{normal}{正规}；
(ii) 对 $A$ 的每个满足 $\dim A_{\mathfrak p}\geq2$ 的\Term{prime-ideal}{素理想} $\mathfrak p$，有 $\operatorname{prof}A_{\mathfrak p}\geq2$。\footnote{参见 EGA IV 5.8.6。}
\end{SGAresult}
\begin{proof}[3.10 的证明]
(i)$\Rightarrow$(ii)。\Term{factorial}{析因}环的\Term{localization}{局部化}仍\Term{factorial}{析因}，故得 (ii) a)。又 $A$ \Term{normal}{正规}，由 $\dim A\geq2$ 及 3.11 (ii) 得 $\operatorname{prof}A\geq2$。最后 $A$ \Term{parafactorial}{旁析因}，因为 $\operatorname{Pic}(X)\simeq\operatorname{Cl}(X')=0$（参见 3.9）。
(ii)$\Rightarrow$(i)。先应用 Serre 判据证明 $A$ \Term{normal}{正规}。因 $\dim A\geq2$，判据的条件 (i) 已含于假设。又对每个 $\mathfrak p\in X$，$A_{\mathfrak p}$ \Term{factorial}{析因}，故\Term{normal}{正规}，因而当 $\dim A_{\mathfrak p}\geq2$ 时\Term{depth}{深度}至少为 $2$。最后，由 (ii) b)，$\operatorname{prof}A\geq2$。再应用 3.9 即可。
\end{proof}
概括以上结果：
\begin{SGAresult}{命题}{3.12}
设 $X$ 为\Term{locally-noetherian}{局部诺特}\Term{prescheme}{预概形}，$\mathcal I$ 为 $X$ 的\Term{quasi-coherent}{拟凝聚}\Term{ideal}{理想}。设 $Y=V(\mathcal I)$，$p\in\mathbb N$。假设：
1) 有 $\operatorname{Leff}(X,Y)$（第 X 讲）；
2) 当 $i=1,2$ 且 $n\geq p$ 时，$H^i(X,\mathcal I^{n+1}/\mathcal I^{n+2})=0$；
3) 对 $Y$ 在 $X$ 中的每个\Term{open-neighborhood}{开邻域} $U$ 及每个 $x\in X-U$，环 $\mathcal O_{X,x}$ \Term{parafactorial}{旁析因}。
\SourcePage{113}
则对每个 $n\geq p$ 及 $Y$ 的每个\Term{open-neighborhood}{开邻域} $U$，下图中的\Term{homomorphism}{同态}都是\Term{isomorphism}{同构}：
\[\begin{tikzcd}\operatorname{Pic}(X)\arrow[r]\arrow[dr]&\operatorname{Pic}(X_n)\\&\operatorname{Pic}(U)\arrow[u]\end{tikzcd}\]
\end{SGAresult}
我们已知道一些\Term{parafactorial}{旁析因}环：
\begin{SGAresult}{定理}{3.13}
(i)（Auslander--Buchsbaum）\Term{regular}{正则}\Term{noetherian}{诺特}\Term{local-ring}{局部环}是\Term{factorial}{析因}环（所以当\Term{dimension}{维数}至少为 $2$ 时\Term{parafactorial}{旁析因}）。\footnote{编者注：可与 Gabber 的如下\Term{purity}{纯性}结果比较：设 $X$ 是\Term{dimension}{维数} $3$ 的\Term{regular}{正则}\Term{local-ring}{局部环} $A$ 的谱，$a$ 的\Term{differential}{微分}非零，即 $a\in\mathfrak m-\mathfrak m^2$，$U$ 是 $V(a)$ 的补集。则 $U$ 上的\Term{vector-bundle}{向量丛}是自由的。简单证明见 Swan R.G., ``A simple proof of Gabber's theorem on projective modules over a localized local ring'', Proc. Amer. Math. Soc. 103 (1988), no.~4, pp.~1025--1030。秩 $1$ 的情形是定理 3.13 的特例。关于任意秩\Term{vector-bundle}{向量丛}的\Term{purity}{纯性}结果（解析或\Term{algebra}{代数}情形），见 Gabber O., ``On purity theorems for vector bundles'', Internat. Math. Res. Notices (2002), no.~15, pp.~783--788。}
(ii) \Term{dimension}{维数}至少为 $4$ 且为\Term{complete-intersection}{完全交}的\Term{noetherian}{诺特}\Term{local-ring}{局部环}是\Term{parafactorial}{旁析因}的。
\end{SGAresult}
\begin{SGAresult}{推论}{3.14（Samuel 猜想）}
若\Term{noetherian}{诺特}\Term{local-ring}{局部环} $A$ 为\Term{complete-intersection}{完全交}，且在余维至少为 $3$ 处\Term{factorial}{析因}（即 $\dim A_{\mathfrak p}\leq3$ 蕴涵 $A_{\mathfrak p}$ \Term{factorial}{析因}），则 $A$ \Term{factorial}{析因}。\footnote{编者注：同一思路而更初等的证明，见 Call F. \& Lyubeznik G., ``A simple proof of Grothendieck's theorem on the parafactoriality of local rings'', in Commutative algebra: syzygies, multiplicities, and birational algebra (South Hadley, MA, 1992), Contemp. Math., vol.~159, American Mathematical Society, Providence, RI, 1994, pp.~15--18。}
\end{SGAresult}
\begin{proof}
对 $A$ 的\Term{dimension}{维数}归纳。若 $\dim A\leq3$，由假设 $A$ \Term{factorial}{析因}。若 $\dim A>3$，注意\Term{complete-intersection}{完全交}的\Term{localization}{局部化}仍为\Term{complete-intersection}{完全交}，由归纳假设，$A$ 以外的所有\Term{localization}{局部化}都\Term{factorial}{析因}。由定理 3.13 (ii)，$A$ \Term{parafactorial}{旁析因}，因此由 3.10 \Term{factorial}{析因}。
\end{proof}
\begin{proof}[3.13 (i) 的证明（依 Kaplansky）]
\footnote{这正是 EGA IV 21.11.1 重述的证明。}
设 $A$ 为\Term{regular}{正则}\Term{noetherian}{诺特}\Term{local-ring}{局部环}，置 $\dim A=n$。若 $n=0,1$，结果已知。设 $n\geq2$，对 $n$ 归纳，假设定理对\Term{dimension}{维数}小于 $n$ 的环已证。置 $X'=\operatorname{Spec}(A)$、$X=X'-\{x\}$，其中 $x=\mathfrak r(A)$。
\SourcePage{114}
$A$ 以外的\Term{localization}{局部化}\Term{regular}{正则}且\Term{dimension}{维数}小于 $n$，所以\Term{factorial}{析因}。此外 $\operatorname{prof}A=\dim A\geq2$。所以只须证明 $\operatorname{Pic}(X)=0$（推论 3.10）。设 $L$ 为\Term{invertible}{可逆} $\mathcal O_X$-模，已知可将它延拓为\Term{coherent}{凝聚} $\mathcal O_{X'}$-模 $L'$。存在 $L'$ 的自由 $\mathcal O_X$-模\Term{resolution}{分解}
\[0\longleftarrow L'\longleftarrow L'_1\longleftarrow\cdots\longleftarrow L'_n\longleftarrow0,\]
因为 $A$ 的\Term{cohomology}{上同调}\Term{dimension}{维数}\Term{finite}{有限}。限制到 $X'$，得到\Term{finite}{有限}自由\Term{resolution}{分解}。所以只须证明下列引理。
\end{proof}
\begin{SGAresult}{引理}{3.15}
设 $(X,\mathcal O_X)$ 为\Term{ringed-space}{赋环空间}，$L$ 为\Term{locally-free}{局部自由} $\mathcal O_X$-模，并有由\Term{finite-type}{有限型}\Term{free-module}{自由模}构成的\Term{finite}{有限}\Term{resolution}{分解}。则 $\det(L)\simeq\mathcal O_X$。
\end{SGAresult}
回忆 $\det(L)$ 定义为 $L$ 的最高\Term{exterior-power}{外幂}。在所考虑的情形，$\det(L)\simeq L$，因为 $L$ \Term{invertible}{可逆}，因此引理给出结论。
\begin{proof}
设上述\Term{exact-sequence}{正合列}为
\[0\longleftarrow L_0\longleftarrow L_1\longleftarrow L_2\longleftarrow\cdots\longleftarrow L_n\longleftarrow0,\]
其中 $L_0=L$。因各项都\Term{locally-free}{局部自由}，有
\[\bigotimes_{0\leq i\leq n}(\det(L_i))^{(-1)^i}\simeq\mathcal O_X.\]
所有 $i>0$ 的 $L_i$ 都自由，故其行列式也自由，因此 $L_0=L$ 的行列式也自由。
\end{proof}
还须证明定理的 (ii)。先证明一个可用于归纳的引理。
\begin{SGAresult}{引理}{3.16}
设 $A$ 为\Term{regular}{正则}环的商的\Term{noetherian}{诺特}\Term{local-ring}{局部环}。设 $t\in\mathfrak r(A)$ 为 $A$-\Term{regular}{正则}元，且 $A$ 对 $t$-\Term{adic-topology}{进拓扑}\Term{complete}{完备}。置 $X'=\operatorname{Spec}(A)$、$Y'=V(t)\simeq\operatorname{Spec}(B)$、$B=A/tA$、$X=X'-\{x\}$、$Y=Y'-\{x\}$、$x=\mathfrak r(A)$。假设：
a) 对 $X$ 的每个\Term{closed-point}{闭点} $y$，$\operatorname{prof}\mathcal O_{X,y}\geq2$；
b) $\operatorname{prof}A/tA\geq3$。
则映射 $\operatorname{Pic}(X)\to\operatorname{Pic}(Y)$ 为\Term{injectivity}{单射}。特别地，若 $B$ \Term{parafactorial}{旁析因}，则 $A$ \Term{parafactorial}{旁析因}。
\end{SGAresult}
\begin{proof}
由 X 2.1，a) 蕴涵 $\operatorname{Lef}(X,Y)$。若证明对每个 $n\geq0$ 有 $H^1(Y,P_n)=0$，则由 (2.2) 知 $\operatorname{Pic}(X)\to\operatorname{Pic}(Y)$ \Term{injectivity}{单射}。若此外 $B$ \Term{parafactorial}{旁析因}，则由 3.5 知 $\operatorname{Pic}(Y)=0$，所以 $\operatorname{Pic}(X)=0$；又因 $t$ 为 $A$-\Term{regular}{正则}元，$\operatorname{prof}(A)=3+1\geq2$，由 3.5，$A$ \Term{parafactorial}{旁析因}。设 $\mathcal I=\widetilde{(tA)}$ 为\Term{ideal}{理想} $tA$ \Term{adjunction}{伴随}的 $\mathcal O_{X'}$-模。在第 1 节，对每个 $n\geq0$ 置 $P_n=(\mathcal I^{n+1}/\mathcal I^{n+2})|_Y$。因 $t$ 为 $A$-\Term{regular}{正则}元，$P_n\simeq\mathcal O_Y$。
\SourcePage{115}
所以只须证明 $H^1(Y,\mathcal O_Y)=0$。因 $Y=Y'-\{x\}$ 在 $Y'$ 中开，有\Term{exact-sequence}{正合列}（I (27)）
\[H^1(Y',\mathcal O_{Y'})\longrightarrow H^1(Y,\mathcal O_Y)\longrightarrow H^2_x(Y',\mathcal O_{Y'}),\]
右端项由假设 b) 为零，左端项因 $Y'$ \Term{affine}{仿射}为零。
\end{proof}
\begin{SGAresult}{引理}{3.17}
保留 3.16 的假设，并再假设：
c) 对 $Y$ 的每个\Term{closed-point}{闭点} $y$，有 $\operatorname{prof}\mathcal O_{X,y}\geq3$；
d) $\operatorname{prof}A/tA\geq4$（强于 b)）；
e) 对 $X$ 的每个满足 $y\in Y$ 的\Term{closed-point}{闭点} $y$，环 $\mathcal O_{X,y}$ \Term{parafactorial}{旁析因}。
则映射 $\operatorname{Pic}(X)\to\operatorname{Pic}(Y)$ 为\Term{isomorphism}{同构}，特别地，$A$ \Term{parafactorial}{旁析因}当且仅当 $B$ \Term{parafactorial}{旁析因}。
\end{SGAresult}
\begin{proof}
已知（X 2.1）a)、c) 蕴涵 $\operatorname{Leff}(X,Y)$。此外，由刚才的论证，d) 蕴涵对每个 $n\geq0$ 及 $i=1,2$ 有 $H^1(Y,P_n)=0$。又 $Y$ 在 $X$ 中每个\Term{open-neighborhood}{开邻域} $U$ 的补集都由\Term{finite}{有限}个\Term{closed-point}{闭点}组成。由 e) 及定理 3.12，得 $\operatorname{Pic}(X)\to\operatorname{Pic}(Y)$ 为\Term{isomorphism}{同构}。另一方面，$\operatorname{prof}A\geq\operatorname{prof}B\geq2$；由判据 3.5 即得 $A$ \Term{parafactorial}{旁析因}当且仅当 $B$ \Term{parafactorial}{旁析因}。
\end{proof}
\begin{proof}[3.13 (ii) 的证明]
设 $R$ 为\Term{regular}{正则}\Term{noetherian}{诺特}\Term{local-ring}{局部环}，$(t_1,\ldots,t_k)$ 为 $R$-序列。置 $B=R/(t_1,\ldots,t_k)$，并假设 $\dim B\geq4$。须证明 $B$ \Term{parafactorial}{旁析因}。对 $k$ 归纳。若 $k=0$，$B$ \Term{regular}{正则}，故由 3.13 (i) \Term{factorial}{析因}，从而由 3.10 \Term{parafactorial}{旁析因}。设 $k\geq1$，并假设定理对 $k'<k$ 已证。置 $A=R/(t_1,\ldots,t_{k-1})$，则 $B=A/t_kA$。由 3.7，可假设 $B$ \Term{complete}{完备}。由归纳假设，$A$ \Term{parafactorial}{旁析因}。证明可应用引理 3.17。已假设 $B$ \Term{complete}{完备}，因而 $A$ 也\Term{complete}{完备}，故 $A$ 对 $t$-\Term{adic-topology}{进拓扑}\Term{complete}{完备}。若 $x\in X$ 且 $x$ 在 $X$ 中闭，则 $A_x$ 是\Term{dimension}{维数}至少为 $4$、且 $k'<k$ 的\Term{complete-intersection}{完全交}。由归纳假设，$A_x$ \Term{parafactorial}{旁析因}，且其\Term{depth}{深度}至少为 $4$。这给出 a)、c)、e)。又 $\dim A\geq5$，故得 d)。
\end{proof}
\begin{SGAresult}{定理}{3.18}
设 $X$ 为\Term{locally-noetherian}{局部诺特}\Term{prescheme}{预概形}，$\mathcal I$ 为 $X$ 的\Term{coherent}{凝聚}\Term{ideal}{理想}\Term{sheaf}{层}。置 $Y=V(\mathcal I)$，设 $n$ 为整数。假设：
(i) 有 $\operatorname{Leff}(X,Y)$（参见例 X 2.1、X 2.2）；
(ii) 对每个 $p\geq n$，当 $i=1,2$ 时有 $H^i(Y,\mathcal I^{p+1}/\mathcal I^{p+2})=0$；
(iii) 对每个\Term{open-subset}{开集} $U\supset Y$ 及每个 $x\in X-U$，环 $\mathcal O_{X,x}$ 或为\Term{dimension}{维数}至少为 $2$ 的\Term{regular}{正则}环，或为\Term{dimension}{维数}至少为 $4$ 的\Term{complete-intersection}{完全交}。
则对每个\Term{open-subset}{开集} $U\supset Y$ 及每个整数 $p\geq n$，\Term{homomorphism}{同态}
\[\operatorname{Pic}(X)\longrightarrow\operatorname{Pic}(U)\longrightarrow\operatorname{Pic}(Y_p)\]
均为\Term{isomorphism}{同构}，其中 $Y_p$ 表示\Term{prescheme}{预概形} $(Y,\mathcal O_X/\mathcal I^{p+1})$。
\end{SGAresult}
只须结合 3.12、3.13。
\SourcePage{116}
